\documentclass[12pt,a4paper,twoside]{article}

%% ---- Pakete ----
\usepackage[utf8]{inputenc}
\usepackage[T1]{fontenc}
\usepackage[ngerman]{babel}
\usepackage{lmodern}
\usepackage{microtype}
\usepackage{geometry}
\usepackage{enumitem}
\setlist[itemize,1]{label=-}

\usepackage{setspace}

%% ---- Grafik & Farben ----
\usepackage{graphicx}
\usepackage{xcolor}
\definecolor{autosarblue}{RGB}{0,80,158}
\definecolor{autosargreen}{RGB}{0,140,69}
\definecolor{autosarorange}{RGB}{220,95,0}
\definecolor{autosargray}{RGB}{90,90,90}
\definecolor{lightblue}{RGB}{173,216,230}
\definecolor{lightyellow}{RGB}{255,255,204}
\definecolor{lightgreen}{RGB}{198,239,206}
\definecolor{lightred}{RGB}{255,199,206}
\definecolor{lightgray}{RGB}{230,230,230}

%% ---- TikZ ----
\usepackage{tikz}
\usetikzlibrary{
  shapes.geometric,
  shapes.misc,
  arrows.meta,
  arrows,
  positioning,
  fit,
  backgrounds,
  decorations.pathreplacing,
  decorations.markings,
  calc,
  matrix,
  patterns,
  shadows,
  automata,
  babel
}

%% ---- Tabellen ----
\usepackage{booktabs}
\usepackage{tabularx}
\usepackage{multirow}
\usepackage{array}

%% ---- Aufzählungen & Listen ----

%% ---- Mathematik ----
\usepackage{amsmath}
\usepackage{amssymb}

%% ---- Listings ----
\usepackage{listings}
\lstset{
  basicstyle=\ttfamily\small,
  breaklines=true,
  frame=single,
  backgroundcolor=\color{lightgray!40},
  keywordstyle=\color{autosarblue}\bfseries,
  commentstyle=\color{autosargreen}\itshape,
  numbers=left,
  numberstyle=\tiny\color{autosargray},
  xleftmargin=2em
}
\lstdefinelanguage{json}{
  basicstyle=\ttfamily\small,
  morestring=[b]",
  morecomment=[l]{//},
  stringstyle=\color{autosargreen!80!black},
  showstringspaces=false,
  breaklines=true
}

%% ---- Hyperlinks ----
\usepackage[
  colorlinks=true,
  linkcolor=autosarblue,
  citecolor=autosarblue,
  urlcolor=autosarorange
]{hyperref}
\usepackage{cleveref}
\crefname{section}{Kapitel}{Kapitel}
\Crefname{section}{Kapitel}{Kapitel}
\crefname{subsection}{Kapitel}{Kapitel}
\Crefname{subsection}{Kapitel}{Kapitel}

\hypersetup{
	colorlinks=true,
	linkcolor=blue!60!black,
	urlcolor=blue!60!black,
	citecolor=blue!60!black
}

\geometry{left=2.5cm, right=2.5cm, top=2.5cm, bottom=2.5cm}

%% ---- Bibliographie ----
\usepackage{csquotes}

%% ---- Glossar ----
\usepackage[acronym,toc]{glossaries}

\newacronym{autosar}{AUTOSAR}{AUTomotive Open System ARchitecture}
\newacronym{ota}{OTA}{Over-the-Air}
\newacronym{fota}{FOTA}{Firmware Over-the-Air}
\newacronym{sota}{SOTA}{Software Over-the-Air}
\newacronym{ecu}{ECU}{Electronic Control Unit}
\newacronym{sg}{SG}{Steuergerät}
\newacronym{zsg}{ZSG}{Zentralsteuergerät}
\newacronym{psg}{PSG}{Peripheriesteuergerät}
\newacronym{cuge}{CUGE}{Central Update Gateway ECU}
\newacronym{ucm}{UCM}{Update and Configuration Management}
\newacronym{rte}{RTE}{Runtime Environment}
\newacronym{bsw}{BSW}{Basic Software}
\newacronym{swc}{SWC}{Software Component}
\newacronym{secoc}{SecOC}{Secure On-Board Communication}
\newacronym{csm}{CSM}{Crypto Service Manager}
\newacronym{hsm}{HSM}{Hardware Security Module}
\newacronym{tls}{TLS}{Transport Layer Security}
\newacronym{pki}{PKI}{Public Key Infrastructure}
\newacronym{can}{CAN}{Controller Area Network}
\newacronym{uds}{UDS}{Unified Diagnostic Services}
\newacronym{someip}{SOME/IP}{Scalable service-Oriented MiddlewarE over IP}
\newacronym{asil}{ASIL}{Automotive Safety Integrity Level}
\newacronym{tara}{TARA}{Threat Analysis and Risk Assessment}
\newacronym{mac}{MAC}{Message Authentication Code}
\newacronym{sum}{SUM}{Software Update Manifest}
\newacronym{ust}{UST}{Update Schedule Table}
\newacronym{dtc}{DTC}{Diagnostic Trouble Code}
\newacronym{hil}{HiL}{Hardware-in-the-Loop}
\newacronym{post}{POST}{Power-On Self-Test}
\newacronym{ap}{AP}{Adaptive Platform}
\newacronym{cp}{CP}{Classic Platform}
\newacronym{api}{API}{Application Programming Interface}
\newacronym{swpa}{SWPA}{Software Package}
\newacronym{sdv}{SDV}{Software-defined Vehicle}
\newacronym{hpc}{HPC}{High-Performance Computer}
\newacronym{soc}{SoC}{System-on-Chip}
\newacronym{aaos}{AAOS}{Android Automotive OS}
\newacronym{evcc}{EVCC}{Electric Vehicle Communication Controller}
\newacronym{ccs}{CCS}{Combined Charging System}
\newacronym{plc}{PLC}{Powerline Communication}
\newacronym{exi}{EXI}{Efficient XML Interchange}
\newacronym{slac}{SLAC}{Signal Level Attenuation Characterization}
\newacronym{sdp}{SDP}{SECC Discovery Protocol}
\newacronym{bms}{BMS}{Battery Management System}
\newacronym{ee}{E/E}{Elektrik/Elektronik}
\newacronym{nat}{NAT}{Network Address Translation}
\newacronym{vndk}{VNDK}{Vendor Native Development Kit}
\newacronym{mcs}{MCS}{Megawatt Charging System}
\newacronym{posix}{POSIX}{Portable Operating System Interface}
\newacronym{ci}{CI}{Continuous Integration}
\newacronym{cd}{CD}{Continuous Deployment}
\newacronym{hal}{HAL}{Hardware Abstraction Layer}
\newacronym{soa}{SoA}{Service-Oriented Architecture}
\newacronym{dds}{DDS}{Data Distribution Service}
\newacronym{v2x}{V2X}{Vehicle-to-Everything}
\newacronym{adas}{ADAS}{Advanced Driver Assistance Systems}

\makeglossaries

\title{
  {\textbf{Zentrales OTA-Update-Management-System im Software-defined Vehicle}\\
  {\Large Architektur, Protokolle und Szenarien für OTA-Update und der Software-defined Vehicle (SDV)-Architektur}}
}
\author{Stephan Epp}
\date{5. März 2026}


\begin{document}


\maketitle

\tableofcontents
\newpage


\section{Einleitung}
\label{sec:einleitung}

Der \textbf{AUTOSAR-Standard} (\textit{AUTomotive Open System ARchitecture-Standard})
bildet seit seiner Einführung die technische Grundlage für die
Softwarearchitektur moderner Kraftfahrzeuge. Parallel dazu vollzieht die
Automobilindustrie gegenwärtig einen der tiefgreifendsten Paradigmenwechsel
ihrer Geschichte: den Ãœbergang zum \textbf{Software-defined Vehicle (SDV)}.
Dieses Konzept – und seine weitreichenden Konsequenzen für OTA-Update-Systeme
– stehen im Mittelpunkt der vorliegenden Arbeit.

\subsection{SDV}
\label{subsec:einl_sdv}

Das \gls{sdv} bezeichnet ein Fahrzeugkonzept, in dem Fahrzeugfunktionen
nicht mehr primär durch spezialisierte, unveränderliche Hardware realisiert
werden, sondern durch Software, die auf einer konsolidierten
Rechenplattform läuft und über den gesamten Fahrzeuglebenszyklus hinweg
aktualisiert, erweitert und angepasst werden kann~\cite{kugele2018sdv}.
Der Begriff wurde maßgeblich durch Tesla geprägt, das erstmals demonstrierte,
dass Fahrzeugeigenschaften wie Reichweite, Beschleunigung oder
Fahrerassistenzfunktionen \emph{nach der Produktion} per drahtlosem
Software-Update verändert werden können. Inzwischen verfolgen nahezu alle
Tier-1-Automobilhersteller und etablierte Zulieferer diesen Ansatz
systematisch~\cite{halder2020secure}.

\paragraph{Worin besteht der Paradigmenwechsel?}
In einem klassischen Fahrzeug des frühen 21.~Jahrhunderts war jede
Fahrzeugfunktion – Motorsteuerung, Bremsenregelung, Klimaautomatik,
Infotainment – in einem dedizierten Steuergerät (\gls{ecu}) verankert.
Funktionsänderungen erforderten entweder einen Werkstattbesuch (Flashen
mittels Diagnosetool über OBD-Schnittstelle) oder sogar einen
Hardware-Tausch. Das \gls{sdv} löst diese enge Hardware-Funktions-Bindung
auf: Software-Bausteine werden von der konkreten Siliziumgeneration
entkoppelt und können auf leistungsstarken, zentralen
\textit{High-Performance Computern} (\glspl{hpc}) konsolidiert werden.
Ein modernes \gls{sdv}-Fahrzeug kommt mit einem Bruchteil der früher
üblichen ECU-Anzahl aus – bei gleichzeitig erheblich größerem
Funktionsumfang~\cite{broy2006challenges}.

\paragraph{Die Rolle der Softwaremenge.}
Moderne Premium-Fahrzeuge enthalten bereits heute mehr als
\textbf{100 Millionen Zeilen Quellcode} – mehr als ein modernes
Kampfflugzeug. Diese Softwaremenge wächst mit zunehmender Autonomisierung
und Vernetzung weiter. Im \gls{sdv} ist diese Software nicht mehr statisch:
Sicherheitsrelevante Patches müssen innerhalb enger Fristen ausgerollt
werden, neue Assistenzfunktionen werden als \textit{Feature-on-Demand}
nach Kauf freigeschaltet, und Diagnosedaten ermöglichen prädiktive
Wartung auf Flottenebene. All dies setzt eine zuverlässige,
fahrzeugweite \textbf{OTA-Update-Infrastruktur} voraus.

\paragraph{E/E-Architektur im SDV: Zonen statt Domänen.}
Der strukturelle Schlüssel des \gls{sdv} ist der Übergang von der
klassischen \emph{domänenbasierten} zu einer \emph{zonenbasierten}
Elektrik/Elektronik-Architektur (\gls{ee}). Während in der domänenbasierten
Architektur Steuergeräte nach Funktion (Antrieb, Fahrwerk, Karosserie,
Infotainment) geclustert wurden, organisiert die Zonenarchitektur die
Fahrzeugelektronik nach \emph{geographischer Lage} im Fahrzeug
(Vorne-Links, Vorne-Rechts, Hinten, Dach). Jede Zone wird durch ein
\emph{Zonensteuergerät} (Zone Controller) koordiniert, das als
intelligenter Hub für alle lokalen Sensoren und Aktuatoren dient.
Die rechenintensiven Aufgaben – Sensorfusion, Fahrfunktionsberechnungen,
Update-Koordination – werden auf zentrale \glspl{hpc} verlagert, die
mit leistungsstarken \glspl{soc} und POSIX-Betriebssystemen
(Linux, \gls{aaos}) betrieben werden. Dieser Architekturwandel
reduziert die Kabelbaummasse erheblich, vereinfacht die
Hardware-Plattformstrategie und schafft die technische Grundlage
dafür, dass ein einziger Update-Eintrittspunkt das gesamte
Fahrzeugnetzwerk versorgen kann.

\paragraph{SDV und AUTOSAR: Chance und Herausforderung.}
AUTOSAR – als der dominierende Softwarestandard der Fahrzeugelektronik –
reagierte auf den SDV-Trend mit der Einführung der
\textbf{AUTOSAR Adaptive Platform} (ab Release 17-03). Diese ermöglicht
dynamische Dienste, serviceorientierte Kommunikation über \gls{someip}
und \gls{dds} sowie ein integriertes Update-and-Configuration-Management
(\gls{ucm}). Gleichzeitig bleibt die bewährte
\textbf{AUTOSAR Classic Platform} für ressourcenbeschränkte Echtzeit-SGs
im Einsatz – ein Fahrzeug enthält heute typischerweise \emph{beide}
Plattformen parallel. Diese Koexistenz heterogener AUTOSAR-Varianten
auf derselben Fahrzeugelektronik stellt OTA-Update-Systeme vor
besondere Herausforderungen: Klassische Mikrocontroller-SGs ohne
POSIX-Betriebssystem müssen genauso zuverlässig aktualisiert werden
wie leistungsstarke Adaptive-Plattform-Knoten, und die Konsistenz
aller Softwareversionen im Fahrzeug muss zu jedem Zeitpunkt
garantiert werden.

\paragraph{Motivation dieser Arbeit.}
Klassische AUTOSAR-Architekturen sehen eine direkte Verbindung zwischen
einem Backend-Server und jedem einzelnen Steuergerät vor, was zu
erhöhtem Kommunikationsaufwand, Synchronisierungsproblemen und
mangelnder Fehlertoleranz führt~\cite{schneider2026ladekomm}. Im \gls{sdv}
erfordert die zonenbasierte \gls{ee}-Architektur jedoch neue
Koordinationsmechanismen, die weit über die bisherigen AUTOSAR-UCM-Konzepte
hinausgehen. Diese Arbeit schlägt daher eine
\emph{hierarchische Erweiterung des AUTOSAR-Standards} vor, in der ein
\emph{Zentralsteuergerät (ZSG)} als dedizierter OTA-Gateway und
Update-Treuhänder fungiert. Das ZSG empfängt Software-Pakete aus dem
Backend, validiert und speichert sie (,,Parken''), und überträgt sie
zeitgestaffelt und zustandsabhängig an \emph{Peripheriesteuergeräte (PSG)} –
darunter Bremsen-SGs, Radar-SGs und Türsteuerungs-SGs. Die Arbeit erweitert
das bestehende AUTOSAR-Schichtenmodell um die notwendigen Dienste,
definiert neue AUTOSAR-Softwarekomponenten und skizziert ein vollständiges
Aktualisierungsszenario mit formalen Zustandsautomaten und
Kommunikationssequenzdiagrammen.

Die Notwendigkeit, Fahrzeugsoftware während der gesamten Fahrzeuglebensdauer
aktuell und sicher zu halten, hat zur Etablierung von
\textit{Firmware-over-the-Air} (FOTA) und
\textit{Software-over-the-Air} (SOTA) als Kerntechnologien des
\gls{sdv} geführt~\cite{nilsson2008key}. Während AUTOSAR Adaptive (R17 ff.)
OTA-Updates prinzipiell adressiert, fehlt ein konsistentes, fahrzeugweites
Koordinationsmodell, das die gesamte Kette vom Backend bis zum individuellen
Steuergerät abdeckt. Dieses Defizit tritt im \gls{sdv} besonders deutlich
hervor, da hier \glspl{hpc}, Zonensteuergeräte und ressourcenbeschränkte
Mikrocontroller als heterogene Plattform koexistieren.

\subsection{Problemstellung}
\label{subsec:problemstellung}

Der Übergang zum \gls{sdv} verschärft strukturelle Schwächen bestehender
OTA-Update-Architekturen in einem Maß, das einen grundlegenden Redesign
erfordert. Die Probleme lassen sich in fünf Kategorien gliedern:

\begin{itemize}[leftmargin=*]
  \item \textbf{Direkte Backend-ECU-Kommunikation ohne Koordination}:
    In klassischen AUTO\-SAR-Architekturen verwaltet jedes Steuergerät
    seine eigene Backend-Verbindung. In einem \gls{sdv}-Fahrzeug mit
    zonenbasierter Architektur und bis zu 60 Steuergeräten führt dies
    zu exzessivem Kommunikationsaufwand, unkontrollierter Bandbreitennutzung
    und einer enorm vergrößerten Angriffsfläche~\cite{halder2020secure}.
    Ressourcenbeschränkte Classic-AUTOSAR-SGs (Bremsen-SG, Radar-SG)
    verfügen zudem weder über einen eigenen Mobilfunk-Stack noch über
    die kryptographische Leistungsfähigkeit, eine direkte TLS-Verbindung
    zum Backend zu unterhalten.

  \item \textbf{Fehlende fahrzeugweite Synchronisierung}:
    Das \gls{sdv} zeichnet sich durch tief verwobene, funktionale
    Abhängigkeiten zwischen Steuergeräten aus: Ein Radar-SG-Update
    setzt eine kompatible Firmware-Version auf dem zentralen
    ADAS-Verarbeitungs-HPC voraus; ein Bremsenregelungs-Update
    darf nur zusammen mit einem kompatiblen Fahrdynamik-SG-Update
    aktiviert werden. Bestehende AUTOSAR-UCM-Konzepte sind primär
    auf das einzelne Steuergerät ausgerichtet und bieten keine
    Mechanismen für atomare, fahrzeugweite Transaktionen.
    Inkonsistente Software-Versionskombinationen können dabei
    sowohl sicherheitskritische Funktionsausfälle als auch
    schwer diagnostizierbare Interoperabilitätsprobleme verursachen.

  \item \textbf{Unzureichende Rollback-Mechanismen}:
    Im \gls{sdv} mit heterogener Plattform – mehrere klassische
    Mikrocontroller-SGs neben leistungsstarken Adaptive-Plattform-Knoten –
    reicht ein Einzel-SG-Rollback nicht aus. Ein fehlerhaftes Update
    eines verbundenen Subsystems kann eine Kettenreaktion auslösen,
    die mehrere Steuergeräte in einen inkonsistenten Zustand versetzt.
    Bestehende Lösungen fehlt ein koordinierter \emph{Fahrzeug-Rollback},
    der alle betroffenen SGs atomar auf ihre Vor-Update-Versionen
    zurückführt.

  \item \textbf{Heterogenität der Plattformen}:
    Das \gls{sdv} kombiniert POSIX-basierte \glspl{hpc} (AUTOSAR Adaptive
    Platform, Linux, \gls{aaos}) mit klassischen Echtzeit-Mikrocontrollern
    (AUTOSAR Classic Platform), die weder über ein Dateisystem noch
    über einen Prozess-Manager verfügen. Ein einheitliches Update-Protokoll
    muss beide Welten bedienen: Das Staging eines Container-Images auf
    einem Adaptive-Dienst unterscheidet sich grundlegend vom Flash-Schreiben
    auf einem Classic-SG über UDS/SOME/IP. Bestehende Standards adressieren
    diese Heterogenität nicht durchgängig.

  \item \textbf{Sicherheitsanforderungen im \gls{sdv}-Kontext}:
    Der erweiterte Connectivity-Footprint des \gls{sdv} – permanente
    Backend-Verbindung, V2X-Kommunikation, Ladeleitungskommunikation – macht
    das Fahrzeug zu einem attraktiveren Angriffsziel als je zuvor.
    ISO~21434 und UN-Regelung~Nr.\,156 (CSMS/SUMS) verpflichten
    Hersteller zu nachweisbaren, end-to-end-gesicherten Update-Ketten.
    Gleichzeitig sind viele ältere Classic-AUTOSAR-SGs nicht für
    \gls{secoc} oder \gls{tls} ausgerüstet und benötigen einen
    sicherheitsaufteilenden Vermittler~\cite{zelle2020security}.
\end{itemize}

\subsection{Ziele der Arbeit}
\label{subsec:ziele}

Aus der in \cref{subsec:problemstellung} beschriebenen Problemlage
ergibt sich ein klares Anforderungsprofil. Diese Arbeit verfolgt
fünf aufeinander aufbauende Ziele:

\begin{enumerate}[leftmargin=*]
  \item \textbf{AUTOSAR-konforme Architekturerweiterung für das \gls{sdv}}:
    Entwurf eines Zentralsteuergeräts (ZSG) als dedizierter OTA-Gateway,
    der die Dual-Plattform-Heterogenität des \gls{sdv} (Classic + Adaptive)
    transparent überbrückt und vollständig in das AUTOSAR-Softwarekomponenten-
    und Dienstmodell integriert ist. Das ZSG soll dabei die
    Leistungsfähigkeit des fahrzeuginternen \gls{hpc} nutzen, ohne
    die Autonomie und Sicherheitszertifizierung der angebundenen
    \glspl{psg} zu kompromittieren.

  \item \textbf{Spezifikation der fahrzeuginternen Kommunikationsprotokolle}:
    Definition der Nachrichten, Sequenzen und Fehlermodi für die
    Kommunikation zwischen ZSG und heterogenen \glspl{psg}, unter
    Berücksichtigung der im \gls{sdv} eingesetzten Bussysteme
    (CAN-FD, Automotive Ethernet, \gls{someip}) sowie der
    kryptographischen Absicherung durch \gls{secoc} und \gls{tls}.

  \item \textbf{Definition eines fahrzeugweiten Zustandsautomaten}:
    Formale Spezifikation des gesamten Update-Lebenszyklus vom
    Backend-Download bis zur abschließenden Validierung, inklusive
    koordinierter Rollback-Sequenz für den Fall, dass ein oder
    mehrere \glspl{psg} die Aktivierung nicht erfolgreich abschließen.
    Der Zustandsautomat muss \gls{asil}-konforme Update-Fenster
    einhalten und den simultanen Betrieb sicherheitskritischer
    Fahrzeugfunktionen während des Update-Prozesses garantieren.

  \item \textbf{Entwicklung eines zeitgestaffelten Update-Szenarios}:
    Ausarbeitung und Analyse eines vollständigen, realistischen
    Multi-SG-Update-Szenarios für ein \gls{sdv}-Fahrzeug mit
    zonenbasierter Architektur, das Bremsen-SG (\gls{asil}-D),
    Radar-SG (\gls{asil}-B) und Tür-SG (\gls{asil}-A) umfasst.
    Die Laufzeitanalyse quantifiziert Overhead durch kryptographische
    Operationen, interne Ãœbertragungszeiten auf den jeweiligen
    Fahrzeugbussen sowie die Gesamtdauer inkl.\,Rollback-Pfad.

  \item \textbf{Bewertung der Sicherheits- und Zuverlässigkeitseigenschaften}:
    Systematische Analyse des vorgeschlagenen Systems gegen die
    Bedrohungsklassen aus \gls{tara} (ISO~21434), Nachweis der
    Konformität mit den strukturellen Anforderungen von
    UN-R\,156 (SUMS) sowie vergleichende Einordnung gegenüber
    bestehenden proprietären und standardisierten OTA-Ansätzen.
\end{enumerate}

\paragraph{Abgrenzung.}
Diese Arbeit konzentriert sich auf die \emph{Architektur} und
\emph{Protokolle} des Update-Systems. Die Implementierung
konkreter AUTOSAR-Softwarekomponenten in C/C++, die Entwicklung
eines HiL-Testaufbaus sowie die formale Verifikation der
Zustandsautomaten sind als explizite Weiterführungsthemen benannt
(\cref{sec:ausblick}) und liegen außerhalb des Rahmens
dieser Arbeit.

\section{Grundlagen}
\label{sec:grundlagen}

\subsection{AUTOSAR Classic und Adaptive}
\label{subsec:autosar}

AUTOSAR (\textit{AUTomotive Open System ARchitecture}) ist ein weltweiter
Entwicklungsstandard für automobile Software, der von einem Konsortium aus
Automobilherstellern, Zulieferern und Werkzeuganbietern entwickelt wurde.
Der Standard gliedert sich in zwei Hauptvarianten:

\begin{description}
  \item[AUTOSAR Classic Platform (CP):] Für ressourcenbeschränkte,
    harte Echtzeit-Steuer\-geräte. Definiert ein statisches Laufzeitsystem
    (RTE), Basissoftware-Module (BSW) und standardisierte Schnittstellen.
  \item[AUTOSAR Adaptive Platform (AP):] Für leistungsstarke,
    POSIX-basierte Steuergeräte, die dynamische Dienste, serviceorientierte
    Kommunikation (SOME/IP) und Update-und-Konfigurations-Management (UCM)
    unterstützen.
\end{description}

\begin{figure}[h!]
  \centering
  %% fig_autosar_layers.tex
%% AUTOSAR Classic Platform - Schichtenmodell (vereinfacht)
\begin{tikzpicture}[
  layer/.style={
    draw=autosarblue, fill=lightblue!60,
    minimum width=10cm, minimum height=0.85cm,
    text centered, font=\small\bfseries
  },
  newlayer/.style={
    draw=autosarblue!80!black, fill=autosarblue!25,
    minimum width=10cm, minimum height=0.85cm,
    text centered, font=\small\bfseries,
    text=autosarblue!80!black
  },
  arrlabel/.style={font=\tiny, text=autosargray}
]

\node[layer] (app)   at (0, 0)    {Applikationsschicht (SWCs)};
\node[layer] (rte)   at (0,-1.0)  {Runtime Environment (RTE)};
\node[newlayer] (otam) at (0,-2.0)
  {{\color{autosarblue}\bfseries OTA-Manager-Dienst (NEU)} – UCM-Erweiterung};
\node[layer] (com)   at (0,-3.0)  {Communication Stack (COM / PDUR / CANIF)};
\node[layer] (diag)  at (0,-4.0)  {Diagnose-Stack (DCM / DEM / UDS)};
\node[layer] (sec)   at (0,-5.0)  {Security-Stack (SecOC / CSM / Crypto)};
\node[layer] (os)    at (0,-6.0)  {Operating System (AUTOSAR OS / POSIX)};
\node[layer] (mcal)  at (0,-7.0)  {Microcontroller Abstraction Layer (MCAL)};
\node[layer] (hw)    at (0,-8.0)  {Hardware (MCU / HSM / Flash / Modem)};

% Verbindungspfeile rechts
\foreach \from/\to in {app/rte, rte/otam, otam/com, com/diag, diag/sec,
                        sec/os, os/mcal, mcal/hw} {
  \draw[-{Stealth[length=3pt]}, autosarblue!50]
    ([xshift=5.2cm]\from.south) -- ([xshift=5.2cm]\to.north);
}

% Legende
\node[draw=autosarblue!80!black, fill=autosarblue!20,
      minimum width=3.2cm, minimum height=0.5cm,
      font=\tiny\bfseries, text=autosarblue!80!black]
  at (4.0,-0.5) {NEU: Erweiterung};
\node[draw=autosarblue, fill=lightblue!60,
      minimum width=3.2cm, minimum height=0.5cm,
      font=\tiny\bfseries]
  at (4.0,-1.2) {Bestehend};

\end{tikzpicture}

  \caption{AUTOSAR-Schichtenmodell (Classic Platform, vereinfacht).
    Die vorgeschlagene Erweiterung fügt den \emph{OTA-Manager-Dienst}
    zwischen Laufzeitsystem und Kommunikationsstack ein.}
  \label{fig:autosar_layers}
\end{figure}

\subsection{Over-the-Air-Updates in der Automobilindustrie}

OTA-Updates bezeichnen das drahtlose Aufspielen von Software-Paketen auf
Fahrzeugsteuergeräte. Man unterscheidet:

\begin{itemize}
  \item \textbf{FOTA (Firmware OTA):} Aktualisierung von Firmware auf
    ressourcenbeschränkten Classic-AUTOSAR-SGs.
  \item \textbf{SOTA (Software OTA):} Aktualisierung von Applikationssoftware
    auf leistungsstärkeren Adaptive-AUTOSAR-SGs.
\end{itemize}

Der Update-Prozess umfasst typischerweise folgende Phasen:
Download $\rightarrow$ Verifikation $\rightarrow$ Staging $\rightarrow$
Installation $\rightarrow$ Aktivierung $\rightarrow$ Validierung.

\subsection{Relevante AUTOSAR-Dienste für OTA}

Der \textit{Update and Configuration Management} (UCM) ist ein Dienst der
AUTOSAR Adaptive Platform, der Software-Pakete verwaltet. Er bietet:
\begin{itemize}
  \item Paketmanagement (Herunterladen, Speichern, Installieren)
  \item Abhängigkeitsprüfung
  \item Rollback-Unterstützung
  \item Authentifizierung über kryptographische Signaturen
\end{itemize}

In der bestehenden AUTOSAR-Architektur ist UCM jedoch primär für ein
einzelnes Steuergerät konzipiert. Eine fahrzeugweite Koordination fehlt.

\begin{figure}[h!]
  \centering
  %% fig_ota_current.tex
%% Konventioneller OTA-Ansatz: Backend -> mehrere ECUs direkt
\begin{tikzpicture}[
  ecustyle/.style={
    draw=autosargray, fill=lightgray!50,
    rounded corners=3pt,
    minimum width=2.6cm, minimum height=0.8cm,
    align=center, font=\small
  },
  backend/.style={
    draw=autosarorange, fill=autosarorange!20,
    rounded corners=4pt,
    minimum width=3.0cm, minimum height=1.2cm,
    align=center, font=\small\bfseries
  },
  arr/.style={-{Stealth[length=5pt]}, thick}
]

\node[backend] (be) at (0,0) {Backend-Server\\{\tiny(OTA-Verwaltung)}};

\node[ecustyle] (bcu)   at (-4.5,-3) {Bremsen-SG\\{\tiny ASIL-D}};
\node[ecustyle] (radar) at (-1.5,-3) {Radar-SG\\{\tiny ASIL-B}};
\node[ecustyle] (door)  at ( 1.5,-3) {Türsteuerungs-SG\\{\tiny QM}};
\node[ecustyle] (ecu4)  at ( 4.5,-3) {Kombi-SG\\{\tiny QM}};

% Verbindungen
\draw[arr, autosarorange!70] (be.south) -- (bcu.north)
  node[midway, left, font=\tiny, xshift=-2pt] {TLS/HTTPS};
\draw[arr, autosarorange!70] (be.south) -- (radar.north)
  node[midway, left=1pt, font=\tiny] {TLS/HTTPS};
\draw[arr, autosarorange!70] (be.south) -- (door.north)
  node[midway, right=1pt, font=\tiny] {TLS/HTTPS};
\draw[arr, autosarorange!70] (be.south) -- (ecu4.north)
  node[midway, right=2pt, font=\tiny] {TLS/HTTPS};

% Problembeschriftung
\node[draw=red!70, fill=red!10, rounded corners=3pt,
      font=\tiny, text width=2.6cm, text centered]
  at (0,-4.8) {\textbf{Probleme:}\\
    Redundante Verbindungen\\Keine Synchronisierung\\
    Gro{\ss}e Angriffsfläche};

\end{tikzpicture}

  \caption{Konventioneller OTA-Ansatz: Jedes Steuergerät kommuniziert
    direkt mit dem Backend-Server (vereinfachte Darstellung).
    Nachteile: hohe Bandbreite, keine Synchronisierung, erhöhte Angriffsfläche.}
  \label{fig:ota_current}
\end{figure}




%% ===========================================================================
\section{Das Software-defined Vehicle (SDV) als Paradigmenwechsel}
\label{sec:sdv}
%% ===========================================================================

\subsection{Definition und Kernkonzept des SDV}
\label{subsec:sdv_definition}

Das \textbf{Software-defined Vehicle (SDV)} bezeichnet ein Fahrzeugkonzept,
bei dem Fahrzeugfunktionen nicht mehr primär durch spezialisierte, unveränderliche
Hardware realisiert werden, sondern durch Software, die auf einer konsolidierten,
leistungsstarken Hardware-Plattform ausgeführt wird und über den gesamten
Fahrzeuglebenszyklus hinweg aktualisiert, erweitert und angepasst werden
kann~\cite{kugele2018sdv}. Der Begriff wurde maßgeblich durch Tesla geprägt,
das erstmals zeigte, dass Fahrzeugeigenschaften wie Beschleunigung, Reichweite
oder Assistenzfunktionen nach der Produktion per Software-Update verändert werden
können. Inzwischen verfolgen nahezu alle Automobilhersteller und Zulieferer
diesen Ansatz~\cite{halder2020secure}.

Die fundamentale Verschiebung des SDV gegenüber klassischen Fahrzeugarchitekturen
lässt sich in drei Dimensionen beschreiben:

\begin{description}
  \item[Hardware-Konsolidierung:] Anstatt für jede Fahrzeugfunktion ein dediziertes
    Steuergerät zu verwenden, werden Funktionen auf wenigen, hochleistungsfähigen
    Rechenplattformen (HPC, SoC) gebündelt. Ein modernes SDV-Fahrzeug kann mit
    wesentlich weniger ECUs auskommen als ein klassisches Fahrzeug, obwohl es mehr
    Funktionen bietet.
  \item[Software-Entkopplung:] Software ist von der spezifischen Hardware-Generation
    weitgehend entkoppelt. Neue Funktionen, Bugfixes und Sicherheitsupdates können
    über den gesamten Fahrzeuglebenszyklus hinweg per OTA-Update eingespielt werden,
    ohne physischen Werkstattbesuch~\cite{nilsson2008key}.
  \item[Serviceorientierung:] Fahrzeugfunktionen werden als Dienste (Services)
    modelliert, die über standardisierte Schnittstellen (z.\,B. SOME/IP, DDS)
    kommunizieren. Dies ermöglicht dynamische Konfiguration, Skalierung und
    den Austausch von Einzelkomponenten~\cite{zelle2020security}.
\end{description}

\subsection{E/E-Architekturtrends: Von der verteilten zur zentralisierten Architektur}
\label{subsec:ee_architektur}

Die Fahrzeug-E/E-Architektur hat in den vergangenen Jahrzehnten eine deutliche
Entwicklung durchlaufen. \Cref{tab:ee_evolution} fasst die wesentlichen
Entwicklungsstufen zusammen.

\begin{table}[htbp]
  \centering
  \caption{Evolutionsstufen der E/E-Architektur im Automobil}
  \label{tab:ee_evolution}
  \begin{tabularx}{\textwidth}{lXXX}
    \toprule
    \textbf{Generation} & \textbf{Architekturprinzip} & \textbf{Typische ECU-Anzahl}
      & \textbf{OTA-Fähigkeit} \\
    \midrule
    1990er & Vollverteilt, domänenlos & 10--30 & Keine \\
    2000er & Domänenbasiert (Antrieb, Fahrwerk \ldots) & 30--80 & Keine / Werkstatt \\
    2010er & Domänenzentralisiert, Gateway & 80--150 & Teilweise (SOTA/FOTA) \\
    2020er & Zonenarchitektur + HPC & 30--60 & Vollständig (SDV) \\
    2030+ & Fahrzeugrechner-Architektur & $<$20 & Nativ + KI-gestützt \\
    \bottomrule
  \end{tabularx}
\end{table}

\subsubsection{Die Zonenarchitektur als strukturelles Rückgrat des SDV}

Die für das SDV charakteristische \textbf{Zonenarchitektur} unterteilt das
Fahrzeug in geographische Zonen (z.\,B. Vorne-Links, Vorne-Rechts, Hinten-Links,
Hinten-Rechts, Dach)~\cite{schneider2026ladekomm}. Jede Zone wird durch ein
\emph{Zonensteuergerät} (engl. \textit{Zone Controller}) verwaltet, das als
intelligenter Hub für alle Sensoren, Aktuatoren und kleineren Steuergeräte
innerhalb seiner Zone dient. Die Zonenarchitektur bietet folgende Vorteile
gegenüber der klassischen domänenbasierten Architektur:

\begin{itemize}
  \item \textbf{Kabelreduktion:} Durch die geographische Organisation entfallen
    lange Kabelwege quer durch das Fahrzeug. Sensoren und Aktuatoren werden lokal
    an das nächste Zonensteuergerät angeschlossen.
  \item \textbf{Leistungskonzentration:} Rechenintensive Aufgaben werden auf
    zentrale HPCs verlagert, die mit leistungsstarken SoCs und POSIX-basierten
    Betriebssystemen (Linux, Android Automotive OS) ausgeführt werden.
  \item \textbf{Update-Vereinfachung:} Software-Updates können zentral auf den
    HPCs und Zonensteuergeräten eingespielt werden, anstatt jede einzelne ECU
    direkt ansprechen zu müssen.
  \item \textbf{Skalierbarkeit:} Neue Fahrzeugvarianten können durch Software-
    Konfiguration der Zonensteuergeräte realisiert werden, ohne neue Hardware zu
    entwickeln.
\end{itemize}

Die folgende TikZ-Darstellung zeigt schematisch die Zonenarchitektur eines
modernen SDV-Fahrzeugs mit zwei zentralen HPCs und vier Zonensteuergeräten:

\begin{figure}[htbp]
  \centering
  \begin{tikzpicture}[
    font=\small,
    hpc/.style={
      rectangle, rounded corners=4pt,
      draw=autosarblue, fill=autosarblue!15,
      minimum width=3.2cm, minimum height=1.1cm,
      align=center, line width=1.2pt
    },
    zone/.style={
      rectangle, rounded corners=3pt,
      draw=autosargreen, fill=autosargreen!12,
      minimum width=2.4cm, minimum height=0.9cm,
      align=center, line width=1pt
    },
    sensor/.style={
      rectangle, rounded corners=2pt,
      draw=autosarorange, fill=autosarorange!10,
      minimum width=1.5cm, minimum height=0.6cm,
      align=center, font=\tiny
    },
    eth/.style={draw=autosarblue, line width=1pt, -Latex},
    can/.style={draw=autosarorange, dashed, line width=0.8pt, -Latex},
    backbone/.style={draw=autosarblue!70, line width=2pt}
  ]
    %% HPCs oben
    \node[hpc] (hpc1) at (-3,0) {HPC 1\\(Fahrdynamik / ADAS)\\{\tiny Linux / AUTOSAR AP}};
    \node[hpc] (hpc2) at ( 3,0) {HPC 2\\(Infotainment / Konnektivität)\\{\tiny Android Automotive OS}};

    %% Backbone
    \draw[backbone] (hpc1.east) -- (hpc2.west)
      node[midway, above, font=\tiny\bfseries, color=autosarblue!70] {Backbone-Ethernet (1\,Gbit/s)};

    %% Zonensteuergeräte
    \node[zone] (z_vl) at (-5,-2.8) {Zone\\Vorne-Links\\{\tiny $\mu$C + CAN-FD}};
    \node[zone] (z_vr) at (-1,-2.8) {Zone\\Vorne-Rechts\\{\tiny $\mu$C + CAN-FD}};
    \node[zone] (z_hl) at ( 1,-2.8) {Zone\\Hinten-Links\\{\tiny $\mu$C + CAN-FD}};
    \node[zone] (z_hr) at ( 5,-2.8) {Zone\\Hinten-Rechts\\{\tiny $\mu$C + CAN-FD}};

    %% Ethernet-Verbindungen HPC -> Zone
    \draw[eth] (hpc1.south) to[out=270,in=90] (z_vl.north);
    \draw[eth] (hpc1.south) to[out=270,in=90] (z_vr.north);
    \draw[eth] (hpc2.south) to[out=270,in=90] (z_hl.north);
    \draw[eth] (hpc2.south) to[out=270,in=90] (z_hr.north);

    %% Sensoren/Aktuatoren je Zone
    \node[sensor] (s_vl1) at (-6.2,-4.5) {Radar\\vorne};
    \node[sensor] (s_vl2) at (-4.8,-4.5) {Bremsen-SG};
    \draw[can] (z_vl.south) -- (s_vl1.north);
    \draw[can] (z_vl.south) -- (s_vl2.north);

    \node[sensor] (s_vr1) at (-1.8,-4.5) {Tür-SG\\links};
    \node[sensor] (s_vr2) at (-0.2,-4.5) {Licht-SG};
    \draw[can] (z_vr.south) -- (s_vr1.north);
    \draw[can] (z_vr.south) -- (s_vr2.north);

    \node[sensor] (s_hl1) at ( 0.2,-4.5) {Tür-SG\\rechts};
    \node[sensor] (s_hl2) at ( 1.8,-4.5) {Airbag-SG};
    \draw[can] (z_hl.south) -- (s_hl1.north);
    \draw[can] (z_hl.south) -- (s_hl2.north);

    \node[sensor] (s_hr1) at ( 4.2,-4.5) {Ladesystem\\(EVCC)};
    \node[sensor] (s_hr2) at ( 5.8,-4.5) {BMS};
    \draw[can] (z_hr.south) -- (s_hr1.north);
    \draw[can] (z_hr.south) -- (s_hr2.north);

    %% Legende
    \begin{scope}[shift={(-6,-5.8)}]
      \draw[eth] (0,0) -- (0.7,0) node[right, font=\tiny] {Ethernet (100/1000\,Mbit/s)};
      \draw[can] (0,-0.35) -- (0.7,-0.35) node[right, font=\tiny] {CAN-FD (2--5\,Mbit/s)};
    \end{scope}
  \end{tikzpicture}
  \caption{Schematische Zonenarchitektur eines SDV-Fahrzeugs mit zwei zentralen
    HPCs und vier Zonensteuergeräten. OTA-Updates werden über die HPCs empfangen
    und über die Zonensteuergeräte an die angebundenen ECUs verteilt.}
  \label{fig:zonenarchitektur}
\end{figure}

\subsection{High-Performance Computer (HPC) als zentrales Element des SDV}
\label{subsec:hpc}

Der High-Performance Computer (HPC) ist das Herzstück des SDV. Er ist typischerweise
als leistungsstarkes System-on-Chip (SoC) ausgeführt und kombiniert:

\begin{itemize}
  \item Mehrere ARM-Cortex-A-Kerne (oft 8--16 Kerne bei 2--3\,GHz) für
    applikationsseitige Aufgaben
  \item Dedizierte Echtzeit-Kerne (ARM-Cortex-R oder -M) für sicherheitskritische
    und zeitkritische Funktionen
  \item Hardware-Beschleuniger für maschinelles Lernen (NPU/GPU) für
    Fahrerassistenz- und Kamerafunktionen
  \item Kryptographische Hardware-Beschleuniger und Hardware Security Module (HSM)
    für Sicherheitsfunktionen einschließlich OTA-Verifikation
  \item Mehrere Netzwerkcontroller für Ethernet (100\,Mbit/s bis 10\,Gbit/s),
    CAN-FD, LIN und Mobilfunk (4G/5G)
\end{itemize}

Betriebssystemseitig kommen auf HPCs im SDV hauptsächlich POSIX-basierte
Systeme zum Einsatz: Linux (z.\,B. Yocto-basiert), Android Automotive OS (AAOS)
oder hypervisor-basierte Ansätze, bei denen ein Echtzeit-Betriebssystem (RTOS)
und ein General-Purpose-OS (GPOS) parallel ausgeführt werden. Dies ermöglicht
die gleichzeitige Erfüllung von Echtzeit-Anforderungen und die Nutzung moderner
Software-Ökosysteme.

\subsubsection{Android Automotive OS (AAOS) im SDV-Kontext}

Android Automotive OS hat sich als führendes Infotainment-Betriebssystem im SDV
etabliert und findet zunehmend auch für sicherheitsnachrüstungsrelevante Dienste
Verwendung~\cite{schneider2026ladekomm}. AAOS ist eine vollständige Android-Variante,
die direkt auf Fahrzeughardware läuft (nicht auf einem angeschlossenen Smartphone).
Für OTA-relevante Dienste ist die \textbf{Partition-Architektur} von AAOS entscheidend:

\begin{itemize}
  \item \textbf{System-Partition:} Enthält das Android-Kernsystem mit den höchsten
    Privilegien. Hier laufen systemnahe Dienste und das Android-Framework. Ist nicht
    für OEM-spezifische Erweiterungen vorgesehen.
  \item \textbf{Vendor-Partition:} Enthält OEM- und Zulieferer-spezifische Treiber
    und Dienste. Hat Zugriff auf alle notwendigen System-APIs und bietet
    Schutzmechanismen gegen unbefugte Änderungen. Dies ist die geeignete Partition
    für fahrzeugspezifische OTA-Dienste.
  \item \textbf{Product-Partition:} Enthält Endbenutzer-Apps und OEM-Anpassungen.
    Wird über das Vendor Interface von der Vendor-Partition getrennt.
\end{itemize}

Für den in dieser Arbeit vorgeschlagenen ZSG-OTA-Manager bedeutet dies: Läuft
der ZSG auf einem HPC mit AAOS, so ist der \texttt{OTA\_UpdateManager\_SWC}
als nativer Dienst in der \textbf{Vendor-Partition} zu implementieren. Er hat
damit vollen Zugriff auf Netzwerk-APIs, Kryptographie-APIs und das Android Keystore
System, während gleichzeitig Isolation von Nutzer-Apps sichergestellt ist.
Der Android Service Manager übernimmt dabei automatisch die Überwachung und
den Neustart des Dienstes im Fehlerfall~\cite{schneider2026ladekomm}.

\subsection{AUTOSAR Adaptive Platform im SDV}
\label{subsec:autosar_adaptive_sdv}

Die AUTOSAR Adaptive Platform (AP) wurde gezielt für den SDV-Kontext entwickelt
und stellt damit das Software-Framework der Wahl für HPC-basierte Steuergeräte dar.
Im Unterschied zur Classic Platform unterstützt die AP:

\begin{itemize}
  \item \textbf{Dynamische Service-Erkennung:} Dienste registrieren sich dynamisch
    und werden zur Laufzeit gefunden (Service Discovery via SOME/IP oder DDS),
    anstatt statisch zur Compile-Zeit konfiguriert zu werden.
  \item \textbf{Flexible Software-Verteilung:} Die AP unterstützt nativ das
    Ausführen von Software in Containern (z.\,B. Docker-ähnliche Isolation), was
    das Deployment und Update einzelner Dienste unabhängig voneinander ermöglicht.
  \item \textbf{Update and Configuration Management (UCM):} Die AP definiert
    den UCM-Dienst als Standardmechanismus für Software-Updates. Im SDV-Kontext
    muss UCM jedoch fahrzeugweit koordiniert werden, was der zentrale Beitrag
    dieser Arbeit adressiert.
  \item \textbf{POSIX-basierte Ausführungsumgebung:} Entwickler können Standard-
    C++-Bibliotheken und Linux-Werkzeuge nutzen, was die Softwareentwicklung
    erheblich vereinfacht im Vergleich zur Classic Platform.
\end{itemize}

\begin{figure}[htbp]
  \centering
  \begin{tikzpicture}[
    font=\small,
    layer/.style={
      rectangle, draw=autosarblue, fill=autosarblue!8,
      minimum width=11cm, minimum height=0.9cm,
      align=center, line width=1pt
    },
    layergreen/.style={
      rectangle, draw=autosargreen, fill=autosargreen!8,
      minimum width=11cm, minimum height=0.9cm,
      align=center, line width=1pt
    },
    layerorange/.style={
      rectangle, draw=autosarorange, fill=autosarorange!8,
      minimum width=11cm, minimum height=0.9cm,
      align=center, line width=1pt
    },
    layergray/.style={
      rectangle, draw=autosargray, fill=lightgray!40,
      minimum width=11cm, minimum height=0.9cm,
      align=center, line width=0.8pt
    },
    comp/.style={
      rectangle, rounded corners=2pt,
      draw=autosarblue!70, fill=autosarblue!20,
      minimum width=2.3cm, minimum height=0.65cm,
      align=center, font=\tiny, line width=0.8pt
    },
    compg/.style={
      rectangle, rounded corners=2pt,
      draw=autosargreen!70, fill=autosargreen!20,
      minimum width=2.3cm, minimum height=0.65cm,
      align=center, font=\tiny, line width=0.8pt
    }
  ]
    %% Schichten von oben nach unten
    \node[layergreen] (app) at (0, 5.4) {\textbf{Applikationsschicht} (Adaptive Applications / SWCs)};
    \node[layer] (ara) at (0, 4.4) {\textbf{ara::com / ara::crypto / ara::diag / ara::exec} (AUTOSAR Runtime for Adaptive)};
    \node[layerorange] (ucm) at (0, 3.4) {\textbf{UCM / Update Manager} (neu: fahrzeugweiter OTA-Koordinator -- diese Arbeit)};

    %% Dienste-Reihe
    \node[comp] (sm)  at (-4.2, 2.4) {State\\Manager};
    \node[comp] (log) at (-1.8, 2.4) {Log \&\\Trace};
    \node[comp] (diag) at (0.3, 2.4) {Diagnostics\\(DM)};
    \node[compg] (nm)  at (2.4, 2.4) {Network\\Management};
    \node[compg] (sec) at (4.5, 2.4) {Crypto\\Service};

    %\draw[draw=autosargray, fill=lightgray!20, line width=0.6pt]
    %  (-5.5, 1.85) rectangle (5.5, 3.0);
    %\node[font=\tiny\itshape, color=autosargray] at (0, 2.85) {Plattformdienste (Functional Clusters)};

    \node[layergray] (os) at (0, 1.3) {\textbf{POSIX-Betriebssystem} (Linux / QNX / AUTOSAR OS)};
    \node[layergray] (hw) at (0, 0.3) {\textbf{Hardware} (SoC -- Cortex-A + Cortex-R + HSM + NPU)};

    %% Pfeile
    %\draw[-Latex, autosarblue] (app.south) -- (ara.north);
    %\draw[-Latex, autosarblue] (ara.south) -- (ucm.north);
    %\draw[-Latex, autosarblue] (ucm.south) -- (os.north);
    %\draw[-Latex, autosarblue] (os.south) -- (hw.north);
  \end{tikzpicture}
  \caption{AUTOSAR Adaptive Platform im SDV-Kontext. Die hervorgehobene
    Schicht zeigt die in dieser Arbeit vorgeschlagene Erweiterung des
    UCM um einen fahrzeugweiten OTA-Koordinator.}
  \label{fig:ap_sdv}
\end{figure}

\subsection{OTA-Update im SDV: Anforderungen und Herausforderungen}
\label{subsec:ota_sdv_anforderungen}

Das SDV stellt gegenüber klassischen Fahrzeugarchitekturen deutlich höhere
Anforderungen an OTA-Update-Systeme. Die folgende Analyse basiert auf den
Erkenntnissen aus dem industriellen Umfeld, insbesondere dem Vorentwicklungsprojekt
von Vector Informatik zur Integration des EVCC in die Zonenarchitektur~\cite{schneider2026ladekomm}.

\subsubsection{Heterogenität der Plattformen}

Im SDV koexistieren fundamental unterschiedliche Plattformen:

\begin{itemize}
  \item \textbf{HPC / SoC (AUTOSAR Adaptive, Linux, AAOS):} Hohe Rechenleistung,
    großer Speicher, schnelle Netzwerkinterfaces. Native OTA-Unterstützung via UCM
    oder Android OTA (A/B-Updates). Bootzeit: mehrere Sekunden bis Minuten.
  \item \textbf{Zonensteuergeräte (AUTOSAR Classic + Adaptive):} Mittlere
    Rechenleistung, eingeschränkter Flash-Speicher, CAN-FD oder Ethernet.
    Bootzeit: unter 500\,ms. Sicherheitskritische Echtzeit-Anforderungen.
  \item \textbf{Periphere Mikrocontroller (AUTOSAR Classic):} Sehr geringe
    Rechenleistung (typisch 32-Bit ARM Cortex-M, 64--512\,MB Flash), CAN oder LIN.
    Bootzeit: unter 100\,ms. Oft ASIL-D-zertifiziert.
\end{itemize}

Der OTA-Update-Mechanismus muss alle drei Plattformtypen unterstützen und dabei
jeweils die spezifischen Einschränkungen hinsichtlich Rechenleistung,
Speicherkapazität und Netzwerkbandbreite berücksichtigen.

\subsubsection{Kryptographische Leistungsanforderungen}

Die im SDV eingesetzte ISO~15118-Ladekommunikation demonstriert exemplarisch
das kryptographische Leistungsproblem: Ein TLS\,1.3-Handshake mit elliptischen
Kurven (ECDSA, P-256) benötigt auf einem klassischen Mikrocontroller bei
80\,MHz Taktfrequenz typisch 2--5\,Sekunden~\cite{schneider2026ladekomm}.
Auf einem HPC mit Hardware-Beschleuniger dauert derselbe Vorgang unter
200\,Millisekunden. Dies ergibt einen Laufzeit-Vorteil von Faktor 8--10,
wie in realen Testszenarien bestätigt wurde.

\begin{table}[htbp]
  \centering
  \caption{Vergleich kryptographischer Laufzeiten: Mikrocontroller vs. HPC im SDV}
  \label{tab:crypto_perf}
  \begin{tabularx}{\textwidth}{lXXX}
    \toprule
    \textbf{Operation} & \textbf{Classic-$\mu$C\newline(80\,MHz, SW)} &
      \textbf{HPC SoC\newline(2\,GHz, HW-Beschl.)} & \textbf{Faktor} \\
    \midrule
    SHA-256 (1\,MB) & $\approx 800\,\text{ms}$ & $<5\,\text{ms}$ & $>160\times$ \\
    ECDSA-P256 Sign & $\approx 1200\,\text{ms}$ & $<2\,\text{ms}$ & $>600\times$ \\
    ECDSA-P256 Verify & $\approx 2400\,\text{ms}$ & $<4\,\text{ms}$ & $>600\times$ \\
    TLS\,1.3 Handshake & $2000$--$5000\,\text{ms}$ & $<200\,\text{ms}$ & $10$--$25\times$ \\
    AES-256-GCM (1\,MB) & $\approx 400\,\text{ms}$ & $<1\,\text{ms}$ & $>400\times$ \\
    \bottomrule
  \end{tabularx}
\end{table}

Diese Leistungsunterschiede unterstreichen die Notwendigkeit, kryptographisch
intensive OTA-Operationen auf HPCs zu verlagern, während sicherheitskritische,
echtzeitfähige Funktionen auf Mikrocontrollern verbleiben. Genau dies ist das
Prinzip des in \Cref{sec:architektur} vorgeschlagenen verteilten ZSG-Konzepts.

\subsubsection{Startzeit und Bootsequenz}

Ein oft übersehener Aspekt im SDV ist die unterschiedliche Bootzeit der
verschiedenen Plattformen. Mikrocontroller starten in Millisekunden und sind
damit sofort einsatzbereit. HPCs hingegen benötigen für den Start von Linux
oder AAOS typisch 5--30\,Sekunden. Dies führt zu einem strukturellen Designproblem:
Das Fahrzeug braucht sofort nach dem Einschalten bestimmte Grundfunktionen
(z.\,B. Ladekommunikation, Bremse), die keinen Verzug dulden, während die
OTA-Kryptographie und komplexe Protokollverarbeitung auf dem HPC noch nicht
verfügbar sind.

Die in diesem Kontext von Vector Informatik beschriebene Lösung~\cite{schneider2026ladekomm}
zeigt exemplarisch: Der Mikrocontroller übernimmt die ersten Schritte der
Ladesequenz (SLAC, SDP) noch vor der Verfügbarkeit des HPC. Sobald der HPC
hochgefahren ist, übernimmt dieser die kryptographisch anspruchsvollen Teile
der Ladekommunikation. Dieses Prinzip der \emph{gestaffelten Ãœbernahme} muss
auch für OTA-Updates gelten: Der Mikrocontroller initiiert den OTA-Handshake,
der HPC führt die Verifikation und Kryptographie durch.

\subsubsection{Fahrzeugweite Konsistenz bei verteilter Software}

Im SDV stellt die Konsistenz der Software-Versionen über alle Steuergeräte
hinweg eine besondere Herausforderung dar. Da Funktionen auf mehrere Steuergeräte
verteilt sind, müssen Software-Updates atomisch auf allen beteiligten Geräten
eingespielt werden. Ein partiell eingespieltes Update – bei dem beispielsweise
der Algorithmus für eine ADAS-Funktion auf dem HPC aktualisiert ist, der
zugehörige Sensor-Treiber auf dem Zonensteuergerät jedoch nicht – kann zu
inkonsistentem und potenziell gefährlichem Verhalten führen~\cite{halder2020secure}.

\subsection{SDV und Cybersicherheit: ISO~21434 im OTA-Kontext}
\label{subsec:sdv_cybersecurity}

Das SDV erweitert die Angriffsfläche eines Fahrzeugs erheblich. Durch die
Konsolidierung auf wenige zentrale HPCs entstehen neue Risiken: Ein kompromittierter
HPC kann theoretisch auf alle angebundenen Steuergeräte zugreifen. Gleichzeitig
bietet das SDV aber auch neue Schutzmaßnahmen durch die Nutzung moderner
Betriebssystem-Sicherheitsmechanismen (SELinux, SecureEnclave, TrustZone).

Die ISO~21434 definiert als Rahmenwerk für Cybersicherheit im Automobil
(\textit{Cybersecurity Engineering}) einen risikobasierten Ansatz (TARA), der
auch für OTA-Update-Systeme im SDV gilt. Besonders relevant sind im SDV-Kontext:

\begin{itemize}
  \item \textbf{NAT-basierte Netzwerksegmentierung:} Zonensteuergeräte nutzen
    NAT-ähnliche Router, die nur erlaubte Protokolle zwischen externen
    Schnittstellen und dem internen Fahrzeugnetz durchlassen~\cite{schneider2026ladekomm}.
    Dies schützt interne ECUs vor direkten Angriffen von außen.
  \item \textbf{Zertifikats-basierte Authentifizierung (PKI):} Im SDV werden
    alle OTA-Pakete durch digitale Signaturen (PKI, ISO~15118) authentifiziert.
    Zertifikate und private Schlüssel werden im HSM des HPC sicher verwahrt.
  \item \textbf{Netzwerk-Monitoring (IDS/IPS):} Der zentrale HPC kann als
    Intrusion Detection System fungieren und abnormale Kommunikationsmuster
    im Fahrzeugnetzwerk erkennen.
\end{itemize}

\subsection{Das SDV-Prinzip der Software-Verteilung: Funktionale Dekomposition}
\label{subsec:sdv_dekomposition}

Ein zentrales Entwurfsprinzip des SDV ist die \emph{funktionale Dekomposition}:
Software-Funktionen werden gemäß ihren nicht-funktionalen Anforderungen
(Echtzeit, Rechenleistung, Sicherheitsniveau, Updatehäufigkeit) auf geeignete
Hardware-Plattformen verteilt. \Cref{fig:sdv_decomposition} illustriert dieses
Prinzip für den OTA-Update-Anwendungsfall.

\begin{figure}[htbp]
  \centering
  \begin{tikzpicture}[
    font=\small,
    box/.style={
      rectangle, rounded corners=4pt,
      draw=#1, fill=#1!12,
      minimum width=4.8cm, minimum height=0.8cm,
      align=center, line width=1pt
    },
    platform/.style={
      rectangle, rounded corners=6pt,
      draw=#1, fill=#1!6,
      inner sep=8pt, line width=1.4pt
    },
    arr/.style={-Latex, draw=autosargray, line width=1pt}
  ]
    %% HPC-Box
    \node[platform=autosarblue, minimum width=5.2cm, minimum height=6.8cm]
      (hpc_box) at (-3.5, 0) {};
    \node[above, font=\small\bfseries, color=autosarblue] at (hpc_box.north)
      {HPC (AUTOSAR Adaptive / AAOS)};

    \node[box=autosarblue] (uc1) at (-3.5,  2.3) {Backend-Kommunikation\\(TLS\,1.3, 5G)};
    \node[box=autosarblue] (uc2) at (-3.5,  1.1) {Paket-Download \&\\Entschlüsselung};
    \node[box=autosarblue] (uc3) at (-3.5, -0.1) {Kryptograph. Verifikation\\(ECDSA, SHA-256)};
    \node[box=autosarblue] (uc4) at (-3.5, -1.3) {Zertifikatsverwaltung\\(PKI, Keystore)};
    \node[box=autosarblue] (uc5) at (-3.5, -2.5) {OTA-Scheduling \&\\Koordination};

    %% Mikrocontroller-Box
    \node[platform=autosargreen, minimum width=5.2cm, minimum height=5.0cm]
      (mc_box) at ( 3.5, -0.5) {};
    \node[above, font=\small\bfseries, color=autosargreen] at (mc_box.north)
      {Mikrocontroller (AUTOSAR Classic)};

    \node[box=autosargreen] (mc1) at (3.5,  0.9) {I/O-Steuerung\\(Hardware-Signale)};
    \node[box=autosargreen] (mc2) at (3.5, -0.3) {Bootloader-Ausführung\\(Secure Boot)};
    \node[box=autosargreen] (mc3) at (3.5, -1.5) {Paket-Empfang\\(via SecOC)};

    %% Ethernet-Verbindung
    \draw[arr, line width=1.4pt, color=autosarblue!60]
      (-0.95, 0) -- (0.95, 0)
      node[midway, above, font=\tiny] {Ethernet}
      node[midway, below, font=\tiny] {(Dual Channel)};

    %% Kriterien-Annotation
    \node[font=\tiny\itshape, color=autosarblue, align=left] at (-3.5,-3.5)
      {Kriterien: Rechenleistung, Kryptographie,\\Zertifikatsspeicher, Update-Logik};
    \node[font=\tiny\itshape, color=autosargreen, align=left] at (3.5,-3.1)
      {Kriterien: Echtzeit, ASIL-D,\\schneller Bootstart, Hardware-I/O};
  \end{tikzpicture}
  \caption{Funktionale Dekomposition der OTA-Funktionen im SDV nach dem
    Prinzip der nicht-funktionalen Anforderungen. Kryptographisch intensive
    Aufgaben verbleiben auf dem HPC; zeitkritische und sicherheitsrelevante
    Funktionen auf dem Mikrocontroller.}
  \label{fig:sdv_decomposition}
\end{figure}

\subsection{Praxisbeispiel: EVCC-Integration in die Zonenarchitektur}
\label{subsec:evcc_praxis}

Als Referenzbeispiel für die SDV-Prinzipien im OTA-Kontext dient die
Integration des Electric Vehicle Communication Controller (EVCC) in die
Zonenarchitektur, wie sie von Tobias Schneider und Michael Vick bei Vector
Informatik untersucht wurde~\cite{schneider2026ladekomm}.

Das EVCC veranschaulicht exemplarisch die Herausforderung, eine Funktion
mit heterogenen Anforderungen in eine SDV-Zonenarchitektur zu integrieren:
Einerseits umfasst es sicherheitskritische Hardware-Steuerung (Control Pilot,
Verriegelungsmotor der Ladedose, CCS-Signale) gemäß ISO~26262; andererseits
erfordert es ressourcenintensive Kryptographie (TLS\,1.3, PKI-Zertifikatsverwaltung)
gemäß ISO~15118-20 und ISO~21434.

Die von Vector untersuchte verteilte Lösung teilt das EVCC auf zwei Steuergeräte auf:

\begin{enumerate}
  \item \textbf{Mikrocontroller-basiertes Gateway} (MICROSAR Classic):
    Ãœbernimmt die sicherheitskritischen Funktionen (Control Pilot, Hardware-I/Os,
    SLAC, SDP, Verriegelungsmotor), profitiert von Echtzeitfähigkeit und schnellem
    Bootstart ($<100$\,ms). Implementiert einen NAT-ähnlichen Router, der nur
    Ladeprotokoll-Nachrichten in das Fahrzeugnetz durchlässt.
  \item \textbf{HPC mit Android Automotive OS} (MICROSAR Charge):
    Ãœbernimmt die rechenintensive Ladekommunikation (TLS\,1.3, EXI-Encoding,
    XML-Security, Zertifikatsverwaltung für Plug\,\&\,Charge). Der Charge-Service
    wird als nativer Dienst in der Vendor-Partition des AAOS implementiert
    und durch den Android Service Manager überwacht.
\end{enumerate}

Das Ergebnis ist bemerkenswert: Die TLS\,1.3-Laufzeit verbessert sich
um den Faktor 8--10 gegenüber einer rein Mikrocontroller-basierten Lösung,
alle ISO-15118-20-Zeitlimits werden eingehalten, und der Speicherbedarf für
Zertifikatsketten stellt auf der SoC-Plattform kein Problem dar~\cite{schneider2026ladekomm}.

\paragraph{Ãœbertragung auf das OTA-Update-System.}
Die Erkenntnisse aus der EVCC-Integration lassen sich direkt auf das
in dieser Arbeit vorgeschlagene ZSG-OTA-System übertragen:

\begin{itemize}
  \item Der \textbf{Mikrocontroller im ZSG} übernimmt den Bootvorgang und
    die ersten Hand\-shake-Schritte des OTA-Protokolls noch vor der Verfügbarkeit
    des HPC. Er stellt sicher, dass das Fahrzeug auch während eines HPC-Neustarts
    (z.\,B. nach einem Update des AAOS) die Grundfunktionen aufrechterhält.
  \item Der \textbf{HPC im ZSG} übernimmt die kryptographische Verifikation
    aller OTA-Pakete, verwahrt Zertifikate und private Schlüssel im Hardware
    Keystore, und koordiniert die zeitgestaffelte Verteilung an die PSGs.
  \item Die \textbf{Dual-Channel-Ethernet-Verbindung} zwischen Mikrocontroller
    und HPC trennt das OTA-Protokoll von der Statussynchronisierung, analog
    zur EVCC-Architek\-tur.
\end{itemize}

\subsection{SDV als Befähiger: Zukünftige OTA-Szenarien}
\label{subsec:sdv_zukunft}

Das SDV eröffnet OTA-Update-Szenarien, die mit klassischen Fahrzeugarchitekturen
nicht realisierbar wären:

\begin{description}
  \item[Differenzielle Funktions-Updates:] Statt vollständiger Firmware-Images
    können im SDV einzelne Software-Dienste (Microservices) unabhängig aktualisiert
    werden. Ein neuer Parkassistent kann als Dienst auf dem HPC eingespielt werden,
    ohne andere Fahrzeugfunktionen zu beeinträchtigen.
  \item[Feature-on-Demand:] Fahrzeugfunktionen können nach dem Kauf freigeschaltet
    werden (z.\,B. erhöhte Motorleistung, erweiterte Assistenzfunktionen). Dies
    erfordert sichere OTA-Mechanismen für Lizenzierung und Aktivierung.
  \item[KI-modellbasierte Updates:] KI-Netzwerkmodelle für Fahrerassistenz
    und Sprachsteuerung können via OTA aktualisiert werden. Die Paketgrößen können
    dabei Gigabytes erreichen, was neue Anforderungen an Bandbreite,
    Speicherverwaltung und Delta-Update-Algorithmen stellt.
  \item[Multi-ECU-Transaktionen:] Ein funktionales Update (z.\,B. neues ABS-Steuer\-gesetz)
    betrifft gleichzeitig HPC, Zonensteuergerät und Bremsen-SG. Im SDV ermöglicht
    der zentrale ZSG-OTA-Manager atomare Multi-ECU-Transaktionen mit vollständigem
    Rollback bei Teilversagen.
  \item[Continuous Deployment (CD) im Fahrzeug:] In Entwicklungsfahrzeugen und
    in der Flottenverwaltung ermöglicht das SDV Continuous-Integration/Continuous-
    Deploy\-ment-Pipelines (CI/CD), bei denen neue Software-Versionen automatisch
    getestet und eingespielt werden – analog zu modernen Cloud-Deployments.
\end{description}

\begin{table}[h]
  \centering
  \caption{SDV-Merkmale und deren Auswirkungen auf das OTA-Update-System}
  \label{tab:sdv_ota_impact}
  \begin{tabularx}{\textwidth}{lXX}
    \toprule
    \textbf{SDV-Merkmal} & \textbf{OTA-Anforderung} & \textbf{Lösungsansatz (diese Arbeit)} \\
    \midrule
    Zonenarchitektur & Hierarchische Update-Verteilung über Zonensteuergeräte
      & ZSG als zentraler Koordinator mit gestaffelter Verteilung \\
    HPC mit POSIX-OS & Native Kryptographie und PKI-Verwaltung
      & OTA-Manager in Vendor-Partition (AAOS) \\
    Heterogene Plattformen & FOTA und SOTA in einem Update-Zyklus
      & Einheitliches SUM-Manifest mit plattformspezifischen Paketen \\
    Schneller Bootstart $\mu$C & OTA-Handshake vor HPC-Verfügbarkeit
      & Gestaffelter Bootstart: $\mu$C übernimmt erste Schritte \\
    Serviceorientierung & Update einzelner Dienste ohne Systemstopp
      & Modulare SWC-Architektur mit unabhängigen Paket-IDs \\
    5G-Konnektivität & Hohe Bandbreite für große Update-Pakete
      & Backend-Download exklusiv über ZSG mit 5G-Modem \\
    ISO~21434 (Cybersicherheit) & Authentifizierung aller Update-Pakete
      & PKI-Signaturen, TLS\,1.3, HSM-basierte Schlüsselverwaltung \\
    Mehrere Ladedosen (LKW/MCS) & Zentrales Update aller EVCC-Instanzen
      & Einmalige Aktualisierung des Charge-Service im HPC \\
    \bottomrule
  \end{tabularx}
\end{table}


\section{Vorgeschlagene Architekturerweiterung}
\label{sec:architektur}


\subsection{Ãœberblick}

Der Kernvorschlag dieser Arbeit ist die Einführung eines
\textbf{Zentralsteuergeräts (ZSG)} – im Englischen \textit{Central Update
Gateway ECU (CUGE)} – das als einziger Kommunikationspartner des
Backend-Servers agiert und alle Update-Operationen im Fahrzeug koordiniert.

\begin{figure}[h!]
  \centering
  %% fig_architecture_overview.tex
%% Neue Architektur: Backend -> ZSG -> PSGs
\begin{tikzpicture}[
  backend/.style={
    draw=autosarorange, fill=autosarorange!20, rounded corners=4pt,
    minimum width=3.2cm, minimum height=1.4cm,
    align=center, font=\small\bfseries
  },
  zsg/.style={
    draw=autosarblue, fill=autosarblue!20, rounded corners=5pt,
    minimum width=3.8cm, minimum height=2.0cm,
    align=center, font=\small\bfseries,
    line width=1.5pt
  },
  psg/.style={
    draw=autosargray, fill=lightgray!40, rounded corners=3pt,
    minimum width=2.5cm, minimum height=1.0cm,
    align=center, font=\small
  },
  asilD/.style={psg, draw=red!60, fill=red!10},
  asilB/.style={psg, draw=autosargreen!80, fill=lightgreen!60},
  asilQ/.style={psg, draw=autosargray, fill=lightgray!40},
  cloud/.style={
    draw=autosarorange!60, fill=autosarorange!10,
    ellipse, minimum width=2.5cm, minimum height=1.2cm,
    align=center, font=\small
  },
  arr/.style={-{Stealth[length=5pt]}, thick},
  intarr/.style={-{Stealth[length=4pt]}, thick, autosarblue!70},
  annotlabel/.style={font=\tiny, text=autosargray, midway, align=center}
]

%% Cloud / Backend
\node[cloud] (cloud) at (0, 5) {5G/4G\\Backend};

%% ZSG
\node[zsg] (zsg) at (0, 2.2) {
  \textbf{ZSG}\\
  {\small\color{autosarblue}(Central Update Gateway)}\\
  {\tiny OTA-Manager | Paket-Store}\\
  {\tiny Transfer-Agent | SecOC}
};

%% PSGs -- Reihe unten
\node[asilD] (bcu)    at (-5.5, -0.5) {Bremsen-SG\\{\tiny ASIL-D}};
\node[asilD] (mcu)    at (-2.8, -0.5) {Motorsteuerung\\{\tiny ASIL-D}};
\node[asilB] (radar)  at ( 0.0, -0.5) {Radar-SG\\{\tiny ASIL-B}};
\node[asilD] (airbag) at ( 2.8, -0.5) {Airbag-SG\\{\tiny ASIL-D}};
\node[asilQ] (door)   at ( 5.5, -0.5) {Türsteuerungs-SG\\{\tiny QM}};

%% PSGs -- zweite Reihe (weitere Geräte)
\node[asilQ] (kombi) at (-2.2, -2.2) {Kombi-SG\\{\tiny QM}};
\node[asilQ] (infotain) at (2.2, -2.2) {Infotainment-SG\\{\tiny QM}};

%% Verbindungen Backend -> ZSG
\draw[arr, autosarorange, line width=1.5pt]
  (cloud) -- (zsg)
  node[annotlabel, right, xshift=3pt]
    {\textbf{TLS 1.3 / HTTPS}\\MQTT-Benachrichtigung};

%% Verbindungen ZSG -> PSGs
\foreach \psg in {bcu, mcu, radar, airbag, door} {
  \draw[intarr] (zsg.south) -- (\psg.north)
    node[annotlabel, font=\tiny\itshape, xshift=-4pt] {CAN-FD\\SecOC+UDS};
}
\foreach \psg in {kombi, infotain} {
  \draw[intarr] (zsg.south) -- (\psg.north)
    node[annotlabel, font=\tiny\itshape] {Eth/SOME-IP\\SecOC};
}

%% Fahrzeug-Umrandung
\begin{scope}[on background layer]
  \node[draw=autosarblue!30, fill=autosarblue!4,
        rounded corners=10pt, line width=1pt,
        fit=(zsg)(bcu)(mcu)(radar)(airbag)(door)(kombi)(infotain),
        inner sep=8pt, label={[font=\small\bfseries,
        text=autosarblue!70]above right:Fahrzeug}]
        (vehicle) {};
\end{scope}

%% Legende
\matrix[draw=autosargray!50, rounded corners=3pt, fill=white,
        nodes={font=\tiny}, row sep=1pt, column sep=3pt,
        inner sep=5pt] at (6.5, 2) {
  \node[asilD, minimum width=0.8cm, minimum height=0.4cm] {}; &
  \node {ASIL-D (sicherheitskritisch)}; \\
  \node[asilB, minimum width=0.8cm, minimum height=0.4cm] {}; &
  \node {ASIL-B (sicherheitsrelevant)}; \\
  \node[asilQ, minimum width=0.8cm, minimum height=0.4cm] {}; &
  \node {QM (nicht sicherheitskritisch)}; \\
  \node[zsg, minimum width=0.8cm, minimum height=0.4cm,
        line width=1pt] {}; &
  \node {ZSG (neu)}; \\
};

\end{tikzpicture}

  \caption{Überblick der vorgeschlagenen Architektur mit Zentralsteuergerät
    (ZSG) und angebundenen Peripheriesteuergeräten (PSGs).
    Das ZSG ist der einzige Kommunikationsknoten zum Backend.}
  \label{fig:architecture_overview}
\end{figure}

\subsection{Zentralsteuergerät (ZSG)}
\label{subsec:zsg}

Das ZSG ist ein dediziertes, leistungsstarkes Steuergerät auf Basis der
\textbf{AUTOSAR Adaptive Platform}. Es übernimmt folgende Funktionen:

\begin{enumerate}
  \item \textbf{Backend-Kommunikation}: Einzige TLS-gesicherte
    Außenverbindung über ein Mobilfunkmodem (4G/5G).
  \item \textbf{Paket-Download}: Herunterladen und Zwischenspeichern
    (\textit{Parken}) von Update-Paketen für alle PSGs.
  \item \textbf{Kryptographische Verifikation}: Signaturprüfung und
    Integritätsprüfung aller Pakete.
  \item \textbf{Update-Koordination}: Zeitgestaffelte Verteilung der
    Pakete an die PSGs gemäß einem konfigurierbaren Update-Plan.
  \item \textbf{Statusüberwachung}: Monitoring aller PSG-Update-Zustände.
  \item \textbf{Rollback-Management}: Koordinierter Rollback im Fehlerfall.
\end{enumerate}

\begin{figure}[h!]
  \centering
  %% fig_zsg_components.tex
%% ZSG-interne Softwarekomponenten
\begin{tikzpicture}[
  compbox/.style={
    draw=autosarblue, fill=autosarblue!15,
    rounded corners=4pt, minimum width=3.8cm, minimum height=1.0cm,
    align=center, font=\small\bfseries
  },
  existbox/.style={
    draw=autosargray!70, fill=lightgray!40,
    rounded corners=4pt, minimum width=3.8cm, minimum height=1.0cm,
    align=center, font=\small
  },
  conn/.style={-{Stealth[length=4pt]}, thick, autosarblue!60},
  dconn/.style={<->, thick, autosarblue!40, dashed}
]

%% -- Neue SWCs (blau) --
\node[compbox] (otam)   at (0, 7)
  {\textbf{OTA\_UpdateManager\_SWC}\\{\tiny Orchestrierung, Schedule, Status}};
\node[compbox] (store)  at (0, 5.2)
  {\textbf{OTA\_PackageStore\_SWC}\\{\tiny Verschl. Ablage, Delta-Mgmt}};
\node[compbox] (xfer)   at (0, 3.4)
  {\textbf{OTA\_TransferAgent\_SWC}\\{\tiny CAN-FD/Eth, SecOC, Flusskontrolle}};
\node[compbox] (verify) at (0, 1.6)
  {\textbf{OTA\_Verifier\_SWC}\\{\tiny SHA-256, RSASSA-PSS, HSM-API}};
\node[compbox] (modem)  at (0,-0.2)
  {\textbf{OTA\_ModemAdapter\_SWC}\\{\tiny TLS 1.3, HTTPS, MQTT}};

%% -- Bestehende AUTOSAR-Dienste (grau) --
\node[existbox] (secoc)  at (5.5, 3.4)
  {SecOC\\{\tiny(AUTOSAR-Basisdienst)}};
\node[existbox] (csm)    at (5.5, 1.6)
  {Crypto Service Manager\\{\tiny(CSM – AUTOSAR)}};
\node[existbox] (nm)     at (5.5, 5.2)
  {Network Management\\{\tiny(NM – AUTOSAR)}};
\node[existbox] (dcm)    at (5.5, 7.0)
  {Diagnostic Manager\\{\tiny(DCM/DEM – AUTOSAR)}};
\node[existbox] (nvm)    at (-5.5, 5.2)
  {NV Memory Manager\\{\tiny(NvM – AUTOSAR)}};
\node[existbox] (wdg)    at (-5.5, 1.6)
  {Watchdog Manager\\{\tiny(WdgM – AUTOSAR)}};
\node[existbox] (os)     at (-5.5, 3.4)
  {AUTOSAR OS\\{\tiny(Echtzeit-OS)}};

%% Verbindungen neue SWCs
\draw[conn] (otam) -- (store)
  node[midway, right, font=\tiny] {Paket-Metadaten};
\draw[conn] (store) -- (xfer)
  node[midway, right, font=\tiny] {Paket-Daten};
\draw[conn] (xfer) -- (verify)
  node[midway, right, font=\tiny] {Verifikation};
\draw[conn] (verify) -- (modem)
  node[midway, right, font=\tiny] {Download-API};

%% Verbindungen zu bestehenden
\draw[dconn] (xfer)   -- (secoc);
\draw[dconn] (verify) -- (csm);
\draw[dconn] (otam)   -- (nm);
\draw[dconn] (otam)   -- (dcm);
\draw[dconn] (store)  -- (nvm);
\draw[dconn] (modem)  -- (wdg)
  node[midway, above, font=\tiny] {};
\draw[dconn] (xfer)   -- (os);

%% Legende
\node[compbox, minimum width=2.0cm, minimum height=0.4cm,
      font=\tiny\bfseries] at (-3.5, -1.4)
  {Neue SWC (Erweiterung)};
\node[existbox, minimum width=2.0cm, minimum height=0.4cm,
      font=\tiny] at (1.0, -1.4)
  {Bestehender AUTOSAR-Dienst};

\end{tikzpicture}

  \caption{Softwarekomponenten des Zentralsteuergeräts (ZSG).
    Neu eingeführte AUTOSAR-Erweiterungen sind blau hervorgehoben.}
  \label{fig:zsg_components}
\end{figure}

\subsection{Peripheriesteuergeräte (PSGs)}
\label{subsec:psgs}

Die PSGs können sowohl auf der AUTOSAR Classic Platform (ressourcenbeschränkte
SGs wie Bremsen-SG, Türsteuerungs-SG) als auch auf der AUTOSAR Adaptive
Platform (leistungsstärkere SGs wie Radar-SG) basieren.

\begin{table}[htbp]
  \centering
  \caption{Übersicht der Peripheriesteuergeräte und ihrer Eigenschaften}
  \label{tab:psgs}
  \begin{tabularx}{\textwidth}{lXXX}
    \toprule
    \textbf{PSG-Typ} & \textbf{AUTOSAR-Plattform} &
    \textbf{Sicherheitsstufe (ASIL)} & \textbf{Update-Fenster} \\
    \midrule
    Bremsen-SG & Classic (ressourcenarm) & ASIL-D & Nur im Stillstand \\
    Radar-SG & Adaptive (leistungsstark) & ASIL-B & Fahrt möglich (degradiert) \\
    Türsteuerungs-SG & Classic (ressourcenarm) & QM & Jederzeit \\
    Motorsteuerung-SG & Classic & ASIL-D & Nur im Stillstand \\
    Kombiinstrument-SG & Adaptive & QM & Jederzeit \\
    Airbag-SG & Classic & ASIL-D & Nur im Stillstand \\
    \bottomrule
  \end{tabularx}
\end{table}

\subsection{AUTOSAR-Softwarekomponenten der Erweiterung}
\label{subsec:swc}

Die vorgeschlagene Erweiterung führt folgende neue AUTOSAR
Softwarekomponenten (SWCs) ein:

\paragraph{OTA\_UpdateManager\_SWC (ZSG)}
Hauptkomponente des ZSG. Verwaltet den gesamten Update-Lebenszyklus:
\begin{itemize}
  \item Empfang und Parsing von \textit{Software Update Manifest} (SUM)-Dateien
  \item Orchestrierung der Download- und Staging-Operationen
  \item Verwaltung des Update-Zeitplans (\textit{Update Schedule Table}, UST)
\end{itemize}

\paragraph{OTA\_PackageStore\_SWC (ZSG)}
Persistenter Speicher für geparkte Update-Pakete:
\begin{itemize}
  \item Verschlüsselte Ablage in einem dedizierten Flash-Bereich
  \item Verwaltung von Paket-Metadaten (Ziel-PSG, Version, Abhängigkeiten)
  \item Delta-Update-Unterstützung zur Bandbreitenreduzierung
\end{itemize}

\paragraph{OTA\_TransferAgent\_SWC (ZSG $\rightarrow$ PSG)}
Überträgt Pakete vom ZSG zum Ziel-PSG über den internen Fahrzeugbus:
\begin{itemize}
  \item Gesicherter Transfer über AUTOSAR SecOC (Secure On-Board Communication)
  \item Flusskontrolle und Fehlerwiederholung
  \item Fortschrittsüberwachung
\end{itemize}

\paragraph{OTA\_LocalInstaller\_SWC (PSG)}
Minimale Client-Komponente auf Classic-AUTO\-SAR-PSGs:
\begin{itemize}
  \item Empfang und Verifikation des Update-Pakets
  \item Aktivierung im sicheren Boot-Prozess
  \item Statusrückmeldung an das ZSG
\end{itemize}

\begin{figure}[h!]
  \centering
  %% fig_swc_interaction.tex
%% AUTOSAR SWC-Interaktionsdiagramm: ZSG <-> PSG
\begin{tikzpicture}[
  portR/.style={
    draw=autosarblue, fill=autosarblue!25,
    minimum width=1.0cm, minimum height=0.5cm,
    font=\tiny, text centered
  },
  portP/.style={
    draw=autosargreen, fill=autosargreen!25,
    minimum width=1.0cm, minimum height=0.5cm,
    font=\tiny, text centered
  },
  swcbox/.style={
    draw=autosarblue, fill=autosarblue!10,
    rounded corners=4pt,
    minimum width=4.5cm, minimum height=3.5cm,
    font=\small\bfseries
  },
  psgbox/.style={
    draw=autosargray, fill=lightgray!30,
    rounded corners=4pt,
    minimum width=4.0cm, minimum height=3.0cm,
    font=\small\bfseries
  },
  rtebox/.style={
    draw=autosarorange, fill=autosarorange!15,
    minimum width=2.0cm, minimum height=0.8cm,
    align=center, font=\small\bfseries, rounded corners=2pt
  },
  conn/.style={-{Stealth[length=4pt]}, thick, autosarblue!70}
]

%% ZSG SWC Box
\node[swcbox, label={[font=\tiny\bfseries, autosarblue]above:ZSG -- AUTOSAR Adaptive SWC}]
  (zsgswc) at (-3.5, 0) {};
\node[font=\small\bfseries, text=autosarblue] at (-3.5, 1.3) {OTA\_UpdateManager};
\node[portP, anchor=east] (port_out1) at (-1.25,  0.7) {PP: PKG};
\node[portP, anchor=east] (port_out2) at (-1.25,  0.0) {PP: CMD};
\node[portR, anchor=east] (port_in1)  at (-1.25, -0.7) {RP: STATUS};

%% RTE
\node[rtebox] (rte) at (0.5, 0) {RTE\\{\tiny SecOC}};

%% PSG SWC Box
\node[psgbox, label={[font=\tiny\bfseries, autosargray]above:PSG -- AUTOSAR Classic SWC}]
  (psgswc) at (4.5, 0) {};
\node[font=\small] at (4.5, 0.9) {OTA\_LocalInstaller};
\node[portR, anchor=west] (port_r1) at (2.75,  0.7) {RP: PKG};
\node[portR, anchor=west] (port_r2) at (2.75,  0.0) {RP: CMD};
\node[portP, anchor=west] (port_p1) at (2.75, -0.7) {PP: STATUS};

%% Verbindungen durch RTE
\draw[conn, autosarblue]   (port_out1.east) -- (rte.west |- port_out1)
  -- (rte.east |- port_r1) -- (port_r1.west)
  node[midway, above, font=\tiny] {Paketübertragung};
\draw[conn, autosargreen!70] (port_out2.east) -- (rte.west |- port_out2)
  -- (rte.east |- port_r2) -- (port_r2.west)
  node[midway, above, font=\tiny] {Befehle};
\draw[conn, autosarorange]  (port_p1.west)  -- (rte.east |- port_p1)
  -- (rte.west |- port_in1) -- (port_in1.east)
  node[midway, below, font=\tiny] {Statusmeldung};

%% Legende Ports
\node[portP, minimum width=1.4cm] at (-3.5,-2.2) {PP: Provided};
\node[portR, minimum width=1.4cm] at (-1.5,-2.2) {RP: Required};
\node[font=\tiny, text=autosargray] at (1.2,-2.2)
  {Ports gem. AUTOSAR-Notation};

\end{tikzpicture}

  \caption{Interaktion der neu eingeführten AUTOSAR-Softwarekomponenten.
    Ports und Verbindungen sind entsprechend der AUTOSAR-Notation dargestellt.}
  \label{fig:swc_interaction}
\end{figure}

\subsection{Integration in das AUTOSAR-Schichtenmodell}
\label{subsec:schichtenmodell}

Die vorgeschlagene Erweiterung integriert sich nahtlos in das bestehende
AUTO\-SAR-Schichtenmodell. \Cref{fig:autosar_extended} zeigt die erweiterte
Architektur.

\begin{figure}[h!]
  \centering
  %% fig_autosar_extended.tex
%% Erweitertes AUTOSAR-Schichtenmodell für das ZSG (AP + CP Hybrid)
\begin{tikzpicture}[
  layerAP/.style={
    draw=autosarblue, fill=autosarblue!18,
    minimum width=11cm, minimum height=0.75cm,
    text centered, font=\footnotesize\bfseries
  },
  layerNew/.style={
    draw=autosarblue!90!black, fill=autosarblue!35,
    minimum width=11cm, minimum height=0.75cm,
    text centered, font=\footnotesize\bfseries,
    text=autosarblue!90!black, line width=1.2pt
  },
  layerCP/.style={
    draw=autosargray, fill=lightgray!40,
    minimum width=11cm, minimum height=0.75cm,
    text centered, font=\footnotesize
  },
  layerHW/.style={
    draw=black!70, fill=black!15,
    minimum width=11cm, minimum height=0.75cm,
    text centered, font=\footnotesize\bfseries
  },
  sidelabel/.style={font=\tiny\itshape, text=autosargray, rotate=90}
]

%% Zeilen von oben nach unten
\node[layerNew] (app)  at (0, 0)
  {Adaptive Applications: OTA\_UpdateManager | OTA\_PackageStore | OTA\_TransferAgent};
\node[layerNew] (ara)  at (0,-0.9)
  {ara::com / ara::crypto / ara::ucm  (NEU: ara::ota Extension API)};
\node[layerAP]  (fws)  at (0,-1.8)
  {Adaptive Foundation Services (NM, Log, PER, UCM-Master)};
\node[layerAP]  (iam)  at (0,-2.7)
  {Identity \& Access Management (IAM) + Intrusion Detection (IDS)};
\node[layerAP]  (com)  at (0,-3.6)
  {SOME/IP Communication | SecOC | TLS-Stack};
\node[layerAP]  (exec) at (0,-4.5)
  {Execution Management (EM) + State Management (SM)};
\node[layerAP]  (os)   at (0,-5.4)
  {POSIX-compatible OS (Linux / QNX Adaptive)};
\node[layerNew] (hsm)  at (0,-6.3)
  {Hardware Security Module (HSM) API  (NEU: OTA-Key-Management)};
\node[layerCP]  (mcal) at (0,-7.2)
  {MCAL: CAN-FD | Eth | SPI | Flash | Crypto-HW};
\node[layerHW]  (hw)   at (0,-8.1)
  {Hardware: SoC (Cortex-A + Cortex-R) | HSM-Chip | Modem (5G) | NOR-Flash};

%% Legende
\node[layerNew, minimum width=3cm, minimum height=0.45cm,
      font=\tiny\bfseries] at (4.5,-9.0) {Neue Erweiterungsschicht};
\node[layerAP,  minimum width=3cm, minimum height=0.45cm,
      font=\tiny] at (0.5,-9.0) {Bestehende AP-Schicht};
\node[layerCP,  minimum width=3cm, minimum height=0.45cm,
      font=\tiny] at (-3.5,-9.0) {Bestehende CP-Schicht};

\end{tikzpicture}

  \caption{Erweitertes AUTOSAR-Schichtenmodell für das ZSG.
    Neue Schichten und Dienste sind blau hervorgehoben.
    Der OTA-Manager ist als Adaptive Application implementiert.}
  \label{fig:autosar_extended}
\end{figure}


\section{Kommunikationsprotokolle}
\label{sec:protokolle}


\subsection{Backend-zu-ZSG-Kommunikation}

Die Kommunikation zwischen Backend-Server und ZSG erfolgt über:

\begin{itemize}
  \item \textbf{Transportschicht}: TLS 1.3 über 4G/5G-Mobilfunk
  \item \textbf{Anwendungsprotokoll}: HTTPS mit MQTT-Benachrichtigungen
  \item \textbf{Paketformat}: AUTOSAR-UCM-konformes Software-Paket (SWPA)
  \item \textbf{Authentifizierung}: PKI mit fahrzeugindividuellem Zertifikat
\end{itemize}

\subsection{ZSG-zu-PSG-Kommunikation}

Die interne Kommunikation zwischen ZSG und PSGs nutzt den fahrzeuginternen
CAN/CAN-FD oder Ethernet-Bus und wird durch folgende Protokolle abgesichert:

\begin{itemize}
  \item \textbf{AUTOSAR SecOC}: Message Authentication Codes (MACs) für
    jede Nachricht
  \item \textbf{UDS (ISO 14229)}: Diagnoseprotokoll für Flash-Operationen
    (Service 0x34 / 0x36 / 0x37)
  \item \textbf{AUTOSAR FlatMap}: Adressierungsschema für Ziel-PSGs
\end{itemize}

\subsection{Nachrichtenformat des OTA-Update-Protokolls}

\Cref{lst:manifest} zeigt ein vereinfachtes Software Update Manifest (SUM)
im JSON-Format, wie es vom Backend an das ZSG übermittelt wird.

\begin{lstlisting}[language=json,caption={Software Update Manifest (SUM) -- vereinfachtes Beispiel},label={lst:manifest}]
{
  "manifest_version": "2.1",
  "vehicle_id": "VIN_WDB123456789",
  "update_bundle_id": "bundle-2026-03-001",
  "packages": [
    {
      "ecu_id": "BrakeControlUnit_01",
      "package_id": "BCU_SW_v4.2.1",
      "target_version": "4.2.1",
      "current_version": "4.1.8",
      "install_window": "STANDSTILL_ONLY",
      "priority": 1,
      "checksum_sha256": "a1b2c3d4...",
      "signature": "RSASSA-PSS_base64...",
      "delta_package": true
    },
    {
      "ecu_id": "RadarSensor_Front_01",
      "package_id": "RADAR_SW_v2.0.0",
      "target_version": "2.0.0",
      "current_version": "1.9.5",
      "install_window": "DEGRADED_MODE_OK",
      "priority": 2,
      "checksum_sha256": "e5f6a7b8...",
      "signature": "RSASSA-PSS_base64...",
      "delta_package": false
    },
    {
      "ecu_id": "DoorControl_FL_01",
      "package_id": "DCU_SW_v1.3.0",
      "target_version": "1.3.0",
      "current_version": "1.2.7",
      "install_window": "ANYTIME",
      "priority": 3,
      "checksum_sha256": "c9d0e1f2...",
      "signature": "RSASSA-PSS_base64...",
      "delta_package": false
    }
  ]
}
\end{lstlisting}


\section{Update-Szenarien}
\label{sec:szenario}


\subsection{Zustandsautomat des ZSG-Update-Managers}
\label{subsec:zustandsautomat}

Der OTA-Update-Manager des ZSG durchläuft während eines Update-Zyklus
definierte Zustände, die in \Cref{fig:statemachine} dargestellt sind.

\begin{figure}[h!]
  \centering
  %% fig_statemachine.tex
%% Zustandsautomat des OTA-Update-Managers (ZSG)
\begin{tikzpicture}[
  ->, >=Stealth, shorten >=1pt,
  auto, node distance=3.0cm,
  state/.style={
    draw=autosarblue, fill=autosarblue!15,
    circle, minimum size=1.8cm,
    align=center, font=\tiny\bfseries,
    text width=1.6cm
  },
  stateOK/.style={
    draw=autosargreen, fill=lightgreen!60,
    circle, minimum size=1.8cm,
    align=center, font=\tiny\bfseries,
    text width=1.6cm
  },
  stateERR/.style={
    draw=red!70, fill=red!15,
    circle, minimum size=1.8cm,
    align=center, font=\tiny\bfseries,
    text width=1.6cm
  },
  stateINIT/.style={
    draw=black, fill=black,
    circle, minimum size=0.3cm
  },
  trans/.style={font=\tiny, text=autosargray, align=center},
  bend angle=25
]

%% Zustände
\node[stateINIT]    (init)    at (0,  0)   {};
\node[state]        (idle)    at (2.5, 0)  {IDLE};
\node[state]        (notify)  at (5.5, 0)  {NOTIFIED};
\node[state]        (dl)      at (8.5, 0)  {DOWN-\\LOADING};
\node[state]        (verify)  at (8.5,-3)  {VERIFYING};
\node[state]        (staged)  at (5.5,-3)  {STAGED\\(Geparkt)};
\node[state]        (sched)   at (2.5,-3)  {SCHEDULING};
\node[state]        (xfer)    at (2.5,-6)  {TRANS-\\FERRING};
\node[stateOK]      (done)    at (5.5,-6)  {ALL DONE};
\node[stateERR]     (err)     at (8.5,-6)  {ERROR\\/ FAIL};
\node[state]        (rb)      at (8.5,-9)  {ROLL-\\BACK};
\node[stateOK]      (rbdone)  at (5.5,-9)  {ROLLBACK\\DONE};

%% Transitionen
\draw (init)   -- (idle)    node[trans, above] {};
\draw (idle)   -- (notify)  node[trans, above] {};
\draw (notify) -- (dl)      node[trans, above] {};
\draw (dl)     -- (verify)  node[trans, right] {};
\draw (verify) -- (staged)  node[trans, above] {};
\draw (staged) -- (sched)   node[trans, above] {};
\draw (sched)  -- (xfer)    node[trans, right] {};
\draw (xfer)   -- (done)    node[trans, above] {};
\draw (done)   edge[bend left=20] (idle)
               node[pos=0.5, right, trans] {};

%% Fehlerpfade
%\draw (dl)     -- (err)     node[trans, right, xshift=2pt] {};
\draw (verify) -- (err)     node[trans, right] {};
%\draw (xfer)   -- (err)     node[trans, above, xshift=-4pt]
%                              {};
\draw (err)    -- (rb)      node[trans, right] {};
\draw (rb)     -- (rbdone)  node[trans, above] {};
\draw (rbdone) edge[bend right=15] (idle)
               node[pos=0.5, below, trans] {};

%% Zurück zu idle bei Benutzerabbruch
\draw (notify) edge[bend left=40] (idle)
               node[pos=0.5, above, trans] {};

\end{tikzpicture}

  \caption{Zustandsautomat des OTA-Update-Managers (ZSG).
    Blau markierte Zustände gehören zum regulären Update-Zyklus,
    grün markierte Zustände sind erfolgreiche Endzustände,
    rot markierte Zustände signalisieren einen Fehler.}
  \label{fig:statemachine}
\end{figure}

\paragraph*{Erläuterung des Zustandsautomaten.}
Abbildung~\ref{fig:statemachine} zeigt den vollständigen Zustandsautomaten
des OTA-Update-Managers im ZSG mit elf Zuständen.
Er unterscheidet einen \emph{Normalzyklus} (blaue und grüne Zustände) von
einem \emph{Fehlerpfad} (roter und grüner Zustand).

\smallskip
\noindent\textbf{Normalzyklus.}
Der Automat startet im \textbf{Initialzustand} (ausgefüllter schwarzer
Kreis) und wechselt sofort in den Ruhezustand \textbf{IDLE}.
IDLE ist der einzige stabile Wartezustand: Der OTA-Manager überwacht
kontinuierlich eingehende Nachrichten des Backend-Servers, ohne aktiv
Ressourcen zu belegen.
Sobald eine Update-Benachrichtigung eintrifft und erfolgreich
authentifiziert wurde, wechselt der Manager in den Zustand
\textbf{NOTIFIED}.
Von NOTIFIED kann ein Benutzerabbruch (z.\,B.\ explizite Ablehnung durch
den Nutzer oder eine Fahrzeugbedingung, die den Download verhindert) den
Automaten unmittelbar nach \textbf{IDLE} zurückführen.
Andernfalls beginnt der Manager im Zustand \textbf{DOWNLOADING} mit dem
Abruf aller Update-Pakete vom Backend-Server; die Pakete werden temporär
im internen Speicher des ZSG abgelegt.
Im anschließenden Zustand \textbf{VERIFYING} prüft der Manager die
kryptographische Signatur jedes Pakets (SHA-256 mit ECDSA); ein
Verifikationsfehler führt in den Fehlerpfad.
Nach bestandener Prüfung werden die Pakete im Zustand
\textbf{STAGED (Geparkt)} dauerhaft im Flash-Speicher des ZSG
gespeichert und warten auf günstige Installationsbedingungen der
Peripheriesteuergeräte.
Der Manager ermittelt dann im Zustand \textbf{SCHEDULING} die optimale
Installationsreihenfolge unter Berücksichtigung der deklarierten
Installationsbedingungen der PSGs (\texttt{ANYTIME},
\texttt{DEGRADED\_MODE}, \texttt{STANDSTILL}) sowie etwaiger
Paketabhängigkeiten.
Im Zustand \textbf{TRANSFERRING} überträgt der Manager die Pakete
sequenziell oder – sofern voneinander unabhängig – parallelisiert über
den gesicherten Fahrzeugbus (CAN-FD\,/\,Ethernet) an die jeweiligen
PSGs und wartet auf deren Installationsbestätigung.
Sobald alle PSGs den Zustand \texttt{INSTALLED\_OK} gemeldet haben,
wechselt der Automat in den Endzustand \textbf{ALL DONE} (grün
markiert): Eine Erfolgsbestätigung wird an den Backend-Server übermittelt,
und der Automat kehrt in den Ruhezustand \textbf{IDLE} zurück.

\smallskip
\noindent\textbf{Fehlerpfad.}
Schlägt der Download (z.\,B.\ Verbindungsabbruch, unvollständiges Paket)
oder die Signaturverifikation fehl, wechselt der Automat in den
Fehlerzustand \textbf{ERROR\,/\,FAIL} (rot markiert).
Dieser Zustand ist kein Endzustand: Der Manager leitet unmittelbar den
Wiederherstellungsprozess ein und wechselt in den Zustand
\textbf{ROLLBACK}.
In ROLLBACK stellt das ZSG den zuletzt bekannten, validen Software-Stand
auf sich selbst und auf allen betroffenen PSGs wieder her;
bereits teilweise übertragene Pakete werden verworfen.
Nach erfolgreichem Abschluss der Wiederherstellung wechselt der Automat
in den Zustand \textbf{ROLLBACK DONE} (grün markiert), der anzeigt, dass
das System sich in einem konsistenten, betriebsbereiten Zustand befindet.
Abschließend kehrt der Automat nach \textbf{IDLE} zurück und ist bereit
für einen erneuten Update-Versuch.

\subsection{Zeitgestaffeltes Update-Szenario}
\label{subsec:zeitgestaffelt}

Das folgende Szenario beschreibt ein reales Fahrzeug-Update-Ereignis mit
drei heterogenen Peripheriesteuergeräten (Bremsen-SG, Radar-SG, Tür-SG).
Das Fahrzeug ist geparkt; der Zündschlüssel ist abgezogen; das ZSG ist
im Standby-Betrieb aktiv.

\subsubsection{Phasen des Szenarios}

\paragraph{Phase 1 – Benachrichtigung und Download (t = 0–15 min)}
Der Backend-Server sendet eine Update-Benachrichtigung an das ZSG.
Das ZSG wechselt in den Zustand \texttt{DOWNLOADING} und lädt alle drei
Pakete herunter. Nach erfolgreicher Signaturverifikation werden die Pakete
im internen Speicher des ZSG \emph{geparkt} (Zustand \texttt{STAGED}).

\paragraph{Phase 2 – Update des Türsteuerungs-SGs (t = 15–18 min)}
Das Türsteuerungs-SG hat die Installationsbedingung \texttt{ANYTIME}
und wird sofort aktualisiert. Das ZSG überträgt das Paket, das PSG
installiert und bestätigt.

\paragraph{Phase 3 – Update des Radar-SGs (t = 18–35 min)}
Das Radar-SG kann im Degraded-Mode aktualisiert werden (Fahrassistenz
temporär deaktiviert). Das ZSG überträgt das größere Paket, das
Radar-SG wechselt in den Degraded-Mode, installiert und aktiviert die
neue Firmware via Secure Boot.

\paragraph{Phase 4 – Update des Bremsen-SGs (t = 35–42 min)}
Das Bremsen-SG (ASIL-D) darf nur im Stillstand aktualisiert werden.
Da das Fahrzeug bereits geparkt ist, ist die Bedingung erfüllt. Das ZSG
überträgt das Delta-Paket; das Bremsen-SG installiert im sicheren
Boot-Modus und führt die Selbsttestsequenz (POST) durch.

\paragraph{Phase 5 – Finalisierung (t = 42–45 min)}
Alle PSGs haben den Zustand \texttt{INSTALLED\_OK} gemeldet. Das ZSG
sendet eine Erfolgsbestätigung an den Backend-Server und wechselt zurück
in den Zustand \texttt{IDLE}.

\begin{figure}[h!]
  \centering
  %% fig_timeline.tex
%% Zeitlicher Ablauf des gestaffelten Update-Szenarios (Gantt-ähnlich)
\begin{tikzpicture}[
  x=0.22cm, y=1.0cm,
  barTuer/.style={fill=autosarorange!70, draw=autosarorange!90, rounded corners=1pt},
  barRadar/.style={fill=autosargreen!70, draw=autosargreen!90, rounded corners=1pt},
  barBremse/.style={fill=autosarblue!70, draw=autosarblue!90, rounded corners=1pt},
  barZSG/.style={fill=lightgray!80, draw=autosargray, rounded corners=1pt},
  barDL/.style={fill=autosarorange!30, draw=autosarorange!60, rounded corners=1pt},
  phaseLabel/.style={font=\tiny\bfseries, text=autosargray},
  timeLabel/.style={font=\tiny, text=black}
]

%% Zeitachse (0..45 Minuten)
\foreach \t in {0,5,10,15,20,25,30,35,40,45} {
  \draw[autosargray!50, thin] (\t, 0.3) -- (\t, -5.0);
  \node[timeLabel] at (\t, 0.55) {t=\t};
}
\node[font=\small\bfseries] at (-6, 0.55) {Zeit (min)};

%% Y-Achsenbeschriftungen
\node[font=\small\bfseries] at (-6, 0)   {ZSG (Download)};
\node[font=\small\bfseries] at (-6, -1)  {ZSG (Staging)};
\node[font=\small\bfseries] at (-6, -2)  {\textcolor{autosarorange}{Tür-SG}};
\node[font=\small\bfseries] at (-6, -3)  {\textcolor{autosargreen}{Radar-SG}};
\node[font=\small\bfseries] at (-6, -4)  {\textcolor{autosarblue}{Bremsen-SG}};

%% Horizontale Rasterlinien
\foreach \y in {0,-1,-2,-3,-4} {
  \draw[autosargray!25, thin] (0, \y) -- (45, \y);
}

%% ZSG Download Phase (0-15 min)
\fill[barDL] (0, -0.35) rectangle (15, 0.35)
  node[midway, font=\tiny\bfseries, text=autosarorange!80!black]
    {Backend-Download (Pakete 1-3)};

%% ZSG Staging (0-15 parallel mit Download, abgeschlossen bei 15)
\fill[barZSG] (0, -1.35) rectangle (15, -0.65)
  node[midway, font=\tiny, text=autosargray!80!black]
    {Verifikation \& Staging};

%% Tür-SG Update: Transfer 15-16 min, Installation 16-18 min
\fill[barTuer] (15, -2.35) rectangle (16, -1.65)
  node[midway, font=\tiny\bfseries, text=white] {Transf.};
\fill[barTuer, fill opacity=0.6] (16, -2.35) rectangle (18, -1.65)
  node[midway, font=\tiny, text=white] {Installation};
%% Bestätigungspfeil
\draw[-{Stealth[length=4pt]}, autosarorange, thick]
  (18, -2.0) -- (18, -0.95)
  node[font=\tiny, right, autosarorange] {OK};

%% Radar-SG Update: Transfer 18-26 min, Installation 26-35 min
\fill[barRadar] (18, -3.35) rectangle (26, -2.65)
  node[midway, font=\tiny\bfseries, text=white] {Paket-Transfer};
\fill[barRadar, fill opacity=0.6] (26, -3.35) rectangle (35, -2.65)
  node[midway, font=\tiny, text=white] {Install. + SecureBoot};
\draw[-{Stealth[length=4pt]}, autosargreen!80, thick]
  (35, -3.0) -- (35, -0.95)
  node[font=\tiny, right, autosargreen!80] {OK};

%% Bremsen-SG Update: Transfer 35-39 min, Installation 39-42 min
\fill[barBremse] (35, -4.35) rectangle (39, -3.65)
  node[midway, font=\tiny\bfseries, text=white] {Delta-Transfer};
\fill[barBremse, fill opacity=0.6] (39, -4.35) rectangle (42, -3.65)
  node[midway, font=\tiny, text=white] {Install.+POST};
\draw[-{Stealth[length=4pt]}, autosarblue!80, thick]
  (42, -4.0) -- (42, -0.95)
  node[font=\tiny, right, autosarblue!80] {OK};

%% Abschluss-Meldung 42-45
\fill[barZSG, fill opacity=0.7] (42, -1.35) rectangle (45, -0.65)
  node[midway, font=\tiny] {Backend-Bestätigung};

%% Phasenmarkierungen oben
\draw[thick, autosargray!50] (0, 0.8) -- (45, 0.8);
\fill[autosarorange!20] (0, 0.8) rectangle (15, 1.3);
\fill[autosarorange!40] (15, 0.8) rectangle (18, 1.3);
\fill[autosargreen!30]  (18, 0.8) rectangle (35, 1.3);
\fill[autosarblue!30]   (35, 0.8) rectangle (42, 1.3);
\fill[lightgray!60]     (42, 0.8) rectangle (45, 1.3);

\node[phaseLabel] at (7.5,  1.05) {Phase 1: Download};
\node[phaseLabel] at (16.5, 1.05) {Ph.2};
\node[phaseLabel] at (26.5, 1.05) {Phase 3: Radar-Update};
\node[phaseLabel] at (38.5, 1.05) {Ph.4};
\node[phaseLabel] at (43.5, 1.05) {Ph.5};

\end{tikzpicture}

  \caption{Zeitlicher Ablauf des gestaffelten Update-Szenarios.
    Farbkodierung entspricht dem PSG-Typ:
    \textcolor{autosarblue}{\textbf{Bremsen-SG}},
    \textcolor{autosargreen}{\textbf{Radar-SG}},
    \textcolor{autosarorange}{\textbf{Tür-SG}}.}
  \label{fig:timeline}
\end{figure}

\subsection{Sequenzdiagramm}
\label{subsec:sequenz}

\Cref{fig:sequence} zeigt das detaillierte Kommunikationssequenzdiagramm
des beschriebenen Szenarios zwischen Backend, ZSG und den drei PSGs.

\begin{figure}[h!]
  \centering
  %% fig_sequence.tex
%% Sequenzdiagramm: Backend -> ZSG -> PSGs (gestaffeltes Update)
\begin{tikzpicture}[
  lifeline/.style={draw=autosargray, thick},
  actor/.style={
    draw=autosarblue, fill=autosarblue!15,
    rounded corners=3pt, minimum width=2.4cm, minimum height=0.65cm,
    text centered, font=\small\bfseries
  },
  actorPSG/.style={
    draw=autosargray, fill=lightgray!40,
    rounded corners=3pt, minimum width=2.4cm, minimum height=0.65cm,
    text centered, font=\small
  },
  msgarr/.style={-{Stealth[length=4pt]}, thick},
  retarr/.style={-{Stealth[length=4pt]}, thick, dashed},
  notebox/.style={
    draw=lightyellow!80!autosargray, fill=lightyellow!60,
    rounded corners=2pt, font=\tiny, text centered
  }
]

%% Koepfe (x-Positionen: 0 / 3.5 / 7.0 / 10.5 / 14.0)
\node[actor]    (hBE)     at (0,    0) {Backend};
\node[actor]    (hZSG)    at (3.5,  0) {ZSG};
\node[actorPSG] (hTuer)   at (7.0,  0) {T\"{u}r-SG};
\node[actorPSG] (hRadar)  at (10.5, 0) {Radar-SG};
\node[actorPSG] (hBremse) at (14.0, 0) {Bremsen-SG};

%% Lebenslinien
\foreach \x in {0, 3.5, 7.0, 10.5, 14.0} {
  \draw[lifeline] (\x, -0.35) -- (\x, -18.5);
}

%% Phase 1: Download
\node[notebox, minimum width=14.8cm] at (7.0, -0.8)
  {\textbf{Phase 1} -- Benachrichtigung und Download (t = 0--15 min)};

\draw[msgarr] (0,-1.5) -- (3.5,-1.5)
  node[midway, above, font=\tiny] {MQTT: update\_available()};
\draw[msgarr] (3.5,-2.3) -- (0,-2.3)
  node[midway, above, font=\tiny] {GET /manifest (TLS 1.3)};
\draw[retarr] (0,-3.1) -- (3.5,-3.1)
  node[midway, above, font=\tiny] {200 OK: SUM-Manifest};
\draw[msgarr] (3.5,-4.2) -- (0,-4.2)
  node[midway, above, font=\tiny] {GET /pkg/BCU\_v4.2.1 (delta)};
\draw[retarr] (0,-5.0) -- (3.5,-5.0)
  node[midway, above, font=\tiny] {200 OK: 800 kB};
\draw[msgarr] (3.5,-5.7) -- (0,-5.7)
  node[midway, above, font=\tiny] {GET /pkg/RADAR\_v2.0.0};
\draw[retarr] (0,-6.5) -- (3.5,-6.5)
  node[midway, above, font=\tiny] {200 OK: 12 MB};
\draw[msgarr] (3.5,-7.2) -- (0,-7.2)
  node[midway, above, font=\tiny] {GET /pkg/DCU\_v1.3.0};
\draw[retarr] (0,-8.0) -- (3.5,-8.0)
  node[midway, above, font=\tiny] {200 OK: 200 kB};
\draw[autosarblue!50, dashed] (3.35,-8.5) rectangle (3.65,-9.5);
\node[font=\tiny, text=autosarblue, right] at (3.7,-9.0)
  {Signaturpr\"{u}fung (HSM) -- Pakete geparkt};

%% Phase 2: Tuer-SG
\node[notebox, minimum width=14.8cm] at (7.0,-10.2)
  {\textbf{Phase 2} -- T\"{u}r-SG Update (t = 15--18 min, ANYTIME)};

\draw[msgarr] (3.5,-10.8) -- (7.0,-10.8)
  node[midway, above, font=\tiny] {OTA\_CMD: PREPARE\_UPDATE};
\draw[retarr] (7.0,-11.5) -- (3.5,-11.5)
  node[midway, above, font=\tiny] {STATUS: READY};
\draw[msgarr] (3.5,-12.2) -- (7.0,-12.2)
  node[midway, above, font=\tiny] {UDS 0x34: RequestDownload};
\draw[msgarr] (3.5,-13.0) -- (7.0,-13.0)
  node[midway, above, font=\tiny] {UDS 0x36: TransferData (200 kB)};
\draw[msgarr] (3.5,-13.7) -- (7.0,-13.7)
  node[midway, above, font=\tiny] {UDS 0x37: RequestTransferExit};
\draw[retarr] (7.0,-14.4) -- (3.5,-14.4)
  node[midway, above, font=\tiny] {STATUS: INSTALLED\_OK (v1.3.0)};

%% Phase 3: Radar-SG
\node[notebox, minimum width=14.8cm] at (7.0,-15.1)
  {\textbf{Phase 3} -- Radar-SG Update (t = 18--35 min, DEGRADED\_MODE\_OK)};

\draw[msgarr] (3.5,-15.6) -- (10.5,-15.6)
  node[midway, above, font=\tiny] {OTA\_CMD: ENTER\_DEGRADED\_MODE};
\draw[retarr] (10.5,-16.3) -- (3.5,-16.3)
  node[midway, above, font=\tiny] {STATUS: DEGRADED\_ACTIVE};
\draw[msgarr] (3.5,-17.0) -- (10.5,-17.0)
  node[midway, above, font=\tiny] {UDS 0x34/0x36/0x37: 12 MB Transfer};
\draw[retarr] (10.5,-17.7) -- (3.5,-17.7)
  node[midway, above, font=\tiny] {STATUS: INSTALLED\_OK (v2.0.0)};

%% Phasen 4+5 zusammengefasst
\node[notebox, minimum width=14.8cm] at (7.0,-18.2)
  {\textbf{Phase 4+5} -- Bremsen-SG (t=35--42 min) + Backend-Best\"{a}tigung (t=42--45 min)};

\end{tikzpicture}

  \caption{Sequenzdiagramm des zeitgestaffelten OTA-Update-Szenarios.
    Jede Nachricht ist mit dem verwendeten Protokoll annotiert.}
  \label{fig:sequence}
\end{figure}

\subsection{Fehlerbehandlung und Rollback}

Im Falle eines fehlgeschlagenen Updates reagiert das ZSG nach
folgendem Schema:

\begin{enumerate}
  \item Das betroffene PSG meldet \texttt{INSTALL\_FAILED} zurück.
  \item Das ZSG setzt das PSG in den Zustand \texttt{ROLLBACK\_PENDING}.
  \item Die vorherige Firmware-Version (gespeichert im ZSG als Backup)
    wird an das PSG übertragen.
  \item Nach erfolgreichem Rollback wechselt das PSG in den Zustand
    \texttt{ROLLBACK\_OK}.
  \item Das ZSG benachrichtigt den Backend-Server mit einem
    Fehlerbericht (\textit{Diagnostic Trouble Code}, DTC).
\end{enumerate}

\begin{figure}[h!]
  \centering
  %% fig_rollback.tex
%% Rollback-Sequenz bei fehlgeschlagenem Update des Bremsen-SGs
\begin{tikzpicture}[
  lifeline/.style={draw=autosargray, thick},
  actor/.style={
    draw=autosarblue, fill=autosarblue!15,
    rounded corners=3pt, minimum width=2.6cm, minimum height=0.65cm,
    align=center, font=\small\bfseries
  },
  actorPSG/.style={
    draw=red!60, fill=red!10,
    rounded corners=3pt, minimum width=2.6cm, minimum height=0.65cm,
    align=center, font=\small\bfseries
  },
  msgarr/.style={-{Stealth[length=4pt]}, thick},
  retarr/.style={-{Stealth[length=4pt]}, thick, dashed},
  errarr/.style={-{Stealth[length=4pt]}, thick, red!70},
  okarr/.style={-{Stealth[length=4pt]}, thick, autosargreen!80},
  notebox/.style={
    draw=red!40, fill=red!8,
    rounded corners=2pt, font=\tiny, text centered
  },
  okbox/.style={
    draw=autosargreen!50, fill=lightgreen!40,
    rounded corners=2pt, font=\tiny, text centered
  }
]

%% Koepfe (x: 0 / 4.0 / 8.5)
\node[actor]    (hBE)     at (0,    0) {Backend};
\node[actor]    (hZSG)    at (4.0,  0) {ZSG};
\node[actorPSG] (hBremse) at (8.5,  0)
  {Bremsen-SG\\{\tiny (ASIL-D)}};

%% Lebenslinien
\foreach \x in {0, 4.0, 8.5} {
  \draw[lifeline] (\x, -0.35) -- (\x, -13.0);
}

%% Normaler Update-Versuch (schlaegt fehl)
\node[notebox, minimum width=8.8cm] at (4.25,-0.9)
  {Normaler Update-Versuch -- schl\"{a}gt fehl};

\draw[msgarr] (4.0,-1.5) -- (8.5,-1.5)
  node[midway, above, font=\tiny] {OTA\_CMD: PREPARE\_UPDATE};
\draw[retarr] (8.5,-2.2) -- (4.0,-2.2)
  node[midway, above, font=\tiny] {STATUS: READY};
\draw[msgarr] (4.0,-3.0) -- (8.5,-3.0)
  node[midway, above, font=\tiny] {UDS 0x34/0x36: Delta-Paket};
\draw[errarr] (8.5,-4.0) -- (4.0,-4.0)
  node[midway, above, font=\tiny, text=red!70]
    {\textbf{STATUS: INSTALL\_FAILED} (CRC mismatch)};

%% Fehlerbehandlung
\node[notebox, minimum width=8.8cm] at (4.25,-4.8)
  {Rollback-Ausl\"{o}sung durch ZSG};

\draw[errarr] (4.0,-5.4) -- (0,-5.4)
  node[midway, above, font=\tiny, text=red!70]
    {DTC: 0x4210 INSTALL\_FAIL\_BCU};
\draw[msgarr] (4.0,-6.2) -- (8.5,-6.2)
  node[midway, above, font=\tiny] {OTA\_CMD: ROLLBACK\_INIT};
\draw[retarr] (8.5,-7.0) -- (4.0,-7.0)
  node[midway, above, font=\tiny] {STATUS: ROLLBACK\_READY};
\draw[msgarr] (4.0,-7.8) -- (8.5,-7.8)
  node[midway, above, font=\tiny]
    {UDS 0x34/0x36: Backup-Firmware (v4.1.8)};
\draw[msgarr] (4.0,-8.6) -- (8.5,-8.6)
  node[midway, above, font=\tiny] {UDS 0x37: RequestTransferExit};
\draw[retarr] (8.5,-9.4) -- (4.0,-9.4)
  node[midway, above, font=\tiny]
    {STATUS: ROLLBACK\_ACTIVE -- POST l\"{a}uft};

%% Rollback erfolgreich
\node[okbox, minimum width=8.8cm] at (4.25,-10.2)
  {Rollback erfolgreich -- Fahrzeug wieder betriebsbereit};

\draw[okarr] (8.5,-10.8) -- (4.0,-10.8)
  node[midway, above, font=\tiny, text=autosargreen!80]
    {\textbf{STATUS: ROLLBACK\_OK} (v4.1.8 aktiv)};
\draw[okarr] (4.0,-11.6) -- (0,-11.6)
  node[midway, above, font=\tiny, text=autosargreen!80]
    {Rollback-Bericht: BCU rollback\_ok};
\draw[retarr] (0,-12.4) -- (4.0,-12.4)
  node[midway, above, font=\tiny]
    {ACK: Retry-Bundle eingeplant};

\end{tikzpicture}

  \caption{Rollback-Sequenz bei fehlgeschlagenem Update des Bremsen-SGs.}
  \label{fig:rollback}
\end{figure}


\section{Sicherheitsaspekte}
\label{sec:sicherheit}


\subsection{Bedrohungsmodell}

Die Einführung eines zentralen OTA-Gateways konzentriert die Angriffsfläche
auf einen einzigen Knoten. \Cref{tab:threats} listet relevante Bedrohungen
und Gegenmaßnahmen auf.

\begin{table}[htbp]
  \centering
  \caption{Bedrohungsmodell nach TARA (Threat Analysis and Risk Assessment,
    ISO 21434)}
  \label{tab:threats}
  \begin{tabularx}{\textwidth}{lXXl}
    \toprule
    \textbf{Bedrohung} & \textbf{Auswirkung} & \textbf{Gegenmaßnahme} &
    \textbf{Risiko} \\
    \midrule
    Paket-Fälschung & Schadcode auf PSG & PKI-Signatur, SecOC & Niedrig \\
    Man-in-the-Middle & Abgehörte Updates & TLS 1.3, Zertifikats-Pinning
      & Niedrig \\
    Replay-Angriff & Veraltete Firmware & Versionsprüfung, Nonce & Niedrig \\
    ZSG-Kompromittierung & Vollständige Kontrolle & HSM, Secure Boot, IDS
      & Mittel \\
    Denial-of-Service & Update-Blockade & Retry-Logik, Zeitlimit & Mittel \\
    Physischer Zugriff & Direkte Manipulation & Tamper Detection, Gehäuse
      & Gering \\
    \bottomrule
  \end{tabularx}
\end{table}

\subsection{Kryptographische Maßnahmen}

\begin{itemize}
  \item \textbf{Hardware Security Module (HSM)}: Das ZSG enthält ein
    dediziertes HSM für Schlüsselverwaltung und kryptographische Operationen.
  \item \textbf{Secure Boot}: Alle Firmware-Updates werden durch Secure Boot
    verifiziert, bevor sie aktiviert werden.
  \item \textbf{AUTOSAR Crypto Service Manager (CSM)}: Standardisierte
    kryptographische API für Signaturverifikation und Schlüsselverwaltung.
  \item \textbf{SecOC}: Message Authentication Codes schützen die
    internen Busnachrichten zwischen ZSG und PSGs.
\end{itemize}

\subsection{Zuverlässigkeit und Ausfallsicherheit}

\begin{itemize}
  \item \textbf{Watchdog}: Hardware-Watchdog überwacht den ZSG-Update-Prozess.
  \item \textbf{Dual-Bank-Flash}: PSGs mit ASIL-D-Anforderungen nutzen
    Dual-Bank-Flash für atomare Updates.
  \item \textbf{Checksummenverifikation}: CRC32 und SHA-256 auf allen
    Ãœbertragungsschichten.
  \item \textbf{Transaktionales Update}: Alle Update-Schritte sind
    transaktional – bei Abbruch wird automatisch rollback ausgeführt.
\end{itemize}


\section{Evaluation}
\label{sec:evaluation}


\subsection{Vergleich mit bestehenden Ansätzen}

\Cref{tab:comparison} vergleicht den vorgeschlagenen Ansatz mit
ausgewählten bestehenden OTA-Lösungen für AUTOSAR-Fahrzeuge.

\begin{table}[htbp]
  \centering
  \caption{Vergleich OTA-Ansätze für AUTOSAR-Fahrzeuge}
  \label{tab:comparison}
  \begin{tabularx}{\textwidth}{lXXXX}
    \toprule
    \textbf{Merkmal} & \textbf{AUTOSAR UCM Standard} &
    \textbf{Tesla OTA} & \textbf{BMW ConnectedDrive} &
    \textbf{Dieser Ansatz} \\
    \midrule
    Zentrale Koordination & Nein & Ja (proprietär) & Ja (proprietär) & Ja (AUTOSAR-std.) \\
    AUTOSAR-konform & Ja (teil.) & Nein & Nein & Ja (vollst.) \\
    Offline-Staging & Nein & Ja & Ja & Ja \\
    Rollback & Einzel-SG & Fahrzeugweit & Einzel-SG & Fahrzeugweit \\
    Interop. Classic + Adaptive & Nein & Entf. & Teilweise & Ja \\
    Offene Spezifikation & Ja & Nein & Nein & Ja \\
    \bottomrule
  \end{tabularx}
\end{table}

\subsection{Laufzeitanalyse}

Eine analytische Abschätzung der Datenmengen für ein typisches
Fahrzeug-Update-Bundle (3 PSGs, Delta-Pakete):

\begin{itemize}
  \item Bremsen-SG Delta-Paket: $\approx 800\,\text{kB}$
  \item Radar-SG Vollpaket: $\approx 12\,\text{MB}$
  \item Tür-SG Vollpaket: $\approx 200\,\text{kB}$
  \item Gesamtdownload-Zeit (LTE, 10 Mbit/s): $\approx 10{,}4\,\text{s}$
  \item Interne Ãœbertragung CAN-FD (2 Mbit/s): $\approx 52{,}8\,\text{s}$
\end{itemize}


\section{Schluss}

Der vorgeschlagene Ansatz hat folgende Einschränkungen:
\begin{itemize}
  \item Das ZSG ist ein \emph{Single Point of Failure}. Ein redundantes
    Design (primäres und sekundäres ZSG) würde die Verfügbarkeit erhöhen.
  \item Die Update-Fenster für ASIL-D-Geräte schränken die Flexibilität ein.
  \item Sehr alte Classic-AUTOSAR-SGs ohne SecOC-Unterstützung sind nicht
    direkt kompatibel und benötigen einen Protokoll-Adapter.
\end{itemize}

\subsection{Zukünftige Arbeiten}

\begin{itemize}
  \item Integration von \textbf{AUTOSAR AP UCM Master} als Basis für die
    ZSG-Implementie\-rung.
  \item Erprobung in einem Hardware-in-the-Loop (HiL)-Testaufbau mit
    realen AUTOSAR-SGs in einer vollständigen Zonenarchitektur.
  \item Erweiterung um \textbf{KI-gestütztes Scheduling} für das
    Update-Zeitfenster (z.\,B. Prüfung von Fahrzeugzustand, Standort,
    Netzwerkqualität, Batteriefüllstand).
  \item Formale Verifikation der Zustandsautomaten mit UPPAAL oder
    ähnlichen Werkzeugen.
  \item \textbf{SDV-spezifische Erweiterungen}: Integration des ZSG-Konzepts
    in eine vollständige SDV-Zonenarchitektur mit mehreren HPCs und
    Zonensteuergeräten, analog zum EVCC-Ansatz von Vector~\cite{schneider2026ladekomm}.
  \item \textbf{Feature-on-Demand-Unterstützung}: Erweiterung des Software
    Update Manifest (SUM) um Lizenzierungs- und Aktivierungsmetadaten für
    die sichere Freischaltung kostenpflichtiger Fahrzeugfunktionen.
  \item \textbf{Containerisierung}: Untersuchung von OCI-Container-basierten
    Updates für AUTOSAR Adaptive-Dienste auf AAOS-Plattformen, analog zu
    modernen Cloud-Deployment-Praktiken.
  \item \textbf{Multi-Ladedosen-Szenarien}: Anpassung des ZSG-Koordinators
    für elektrisch betriebene LKW mit mehreren EVCC-Instanzen (CCS + MCS),
    bei denen das Ladeprotokoll zentral im HPC verwaltet wird.
\end{itemize}


\subsection{Fazit}
\label{sec:fazit}


Diese Arbeit hat eine AUTOSAR-konforme Erweiterung für das fahrzeugweite
OTA-Update-Management vorgestellt und diese in den Kontext des
\textbf{Software-defined Vehicle (SDV)} eingebettet. Das SDV stellt als
zentrales Leitbild der modernen Fahrzeugentwicklung neue, tiefgreifende
Anforderungen an OTA-Update-Systeme: heterogene Plattformen, kryptographisch
intensive Sicherheitsmechanismen, zonenbasierte E/E-Architekturen und
fahrzeugweite Konsistenzanforderungen.

Das zentrale Zentralsteuergerät (ZSG) übernimmt alle Update-relevanten
Aufgaben – Download, Staging, Verteilung und Koordination – und entlastet
damit die ressourcenbeschränkten Peripheriesteuergeräte. Im SDV-Kontext
profitiert das ZSG von der Leistungsfähigkeit des HPC: Kryptographisch
intensive Operationen (TLS\,1.3-Handshake, ECDSA-Verifikation) werden um den
Faktor 8--25 schneller ausgeführt als auf klassischen Mikrocontrollern~\cite{schneider2026ladekomm}.
Der Mikrocontroller-Anteil des ZSG stellt dabei sicher, dass der OTA-Prozess
auch während des HPC-Bootvorgangs anlaufen kann.

Die vorgeschlagene Architektur bietet folgende Vorteile:
\begin{itemize}
  \item \textbf{Klare Verantwortlichkeit}: Ein dedizierter Update-Gateway
    reduziert die Komplexität in den PSGs.
  \item \textbf{Reduzierte Angriffsfläche}: Nur das ZSG kommuniziert mit
    dem Backend.
  \item \textbf{Offline-Staging}: Updates können vollständig heruntergeladen
    und verifiziert werden, bevor das Fahrzeug offline geht.
  \item \textbf{ASIL-konforme Verteilung}: Sicherheitskritische SGs werden
    nur in geeigneten Fahrzeugzuständen aktualisiert.
  \item \textbf{Standardkonformität}: Vollständige Kompatibilität mit
    AUTOSAR Classic und Adaptive Platform.
\end{itemize}

Das vorgestellte zeitgestaffelte Szenario zeigt, dass selbst ein komplexes
Multi-SG-Update in unter 45 Minuten – einschließlich Rollback-Fähigkeit –
sicher und zuverlässig durchführbar ist. Die Erweiterung ist als offene
AUTOSAR-Spezifikation konzipiert und kann von allen Konsortiumsmitgliedern
implementiert werden.

\subsection{Ausblick}
\label{sec:ausblick}

Die vorliegende Arbeit legt ein konzeptionelles Fundament für eine
AUTOSAR-konforme, fahrzeugweite OTA-Update-Architektur. Die skizzierten
Konzepte – das Zentralsteuergerät als Update-Gateway, das zeitgestaffelte
Verteilungsprotokoll sowie der hierarchische Zustandsautomat – bilden
gemeinsam einen konsistenten Entwurf, der in einer Vielzahl von Richtungen
erheblich vertieft und erweitert werden kann. Der folgende Ausblick
strukturiert diese Erweiterungsmöglichkeiten entlang von sechs zentralen
Entwicklungsachsen.

\subsubsection{Prototypische Implementierung und experimentelle Validierung}

Die in dieser Arbeit durchgeführte Evaluation basiert auf analytischen
Abschätzungen und konzeptionellen Argumenten. Den nächsten natürlichen
Schritt bildet eine \textbf{prototypische Implementierung} auf realer
Hardware. Als Zielplattform bietet sich ein Aufbau bestehend aus einem
leistungsstarken \gls{hpc}-Board (z.\,B. ein Qualcomm SA8295P oder ein
NVIDIA Orin-Modul) als ZSG-Host, ergänzt durch mehrere
Mikrocontroller-Boards (z.\,B. NXP S32K-Familie oder Renesas RH850) als
simulierte \glspl{psg}, an.

Im Rahmen eines solchen \textbf{Hardware-in-the-Loop (HiL)-Testaufbaus}
würden sämtliche in \cref{sec:evaluation} analytisch ermittelten
Laufzeitwerte empirisch verifiziert. Besonders aufschlussreich wäre
dabei der Vergleich der kryptographischen Ausführungszeiten auf
dem \gls{hpc}-Anteil des ZSG (mit \gls{hsm}-Beschleunigung) gegenüber
einem klassischen Mikrocontroller ohne dedizierte Krypto-Hardware.
Vorliegende Studien deuten auf Beschleunigungsfaktoren von 8 bis 25 hin;
diese Werte müssten für die spezifischen Algorithmen (ECDSA-256,
SHA-3, AES-GCM) und die gewählten Zielplattformen konkret ermittelt werden.

Über die reine Laufzeitanalyse hinaus wäre ein vollständiges
\textbf{Fault-Injection-Test\-regime} erforderlich: Gezielte Unterbrechung
der CAN-FD- und Ethernet-Verbindungen auf verschiedenen Stufen des
Update-Prozesses, Simulation von Spannungsausfällen während der
Flash-Programmierung sowie künstlich herbeigeführte Prüfsummenfehler
in den Software-Paketen würden die Robustheit der Rollback-Mechanismen
unter realistischen Bedingungen nachweisen. Das Ergebnis eines solchen
Testregimes wäre eine quantitativ belegte Aussage über die
\emph{Mean Time to Recovery} (MTTR) nach einem fehlgeschlagenen Update
und über die Vollständigkeit der Wiederherstellung in den definierten
Ausgangszustand.

Schließlich bietet sich eine Integration in bestehende
\textbf{AUTOSAR-Entwicklungsumge\-bungen} an, insbesondere in
Vector Informatik CANoe/CANalyzer oder dSPACE VEOS, um
die entwickelten Protokolle im Kontext vollständiger
Fahrzeugnetzwerk-Simulationen zu erproben und zu debuggen.

\subsubsection{Formale Verifikation und Sicherheitsnachweise}

Die in dieser Arbeit präsentierten Zustandsautomaten und
Sequenzdiagramme sind bewusst in einer für Menschen lesbaren,
semi-formalen Notation gehalten. Für sicherheitskritische
Fahrzeugfunktionen mit \gls{asil}-D-Einstufung ist eine solche
Darstellung als Ausgangspunkt sinnvoll, aber nicht hinreichend.
Künftige Arbeiten sollten die Zustandsautomaten in einer
\textbf{formalen Spezifikationssprache} – etwa in \textit{UPPAAL}
(Timed Automata), \textit{SPIN}/Promela oder \textit{TLA+} –
modellieren und automatisiert verifizieren.

Konkret sollten folgende \textbf{formale Eigenschaften} bewiesen werden:

\begin{itemize}
  \item \textbf{Deadlock-Freiheit}: Der Gesamtautomat aus ZSG-Koordinator
    und allen PSG-Satelliten-Automaten hat keinen erreichbaren
    Deadlock-Zustand, unabhängig von der Anzahl der beteiligten \glspl{psg}
    und der Reihenfolge von Timeout- und Fehlerereignissen.
  \item \textbf{Terminierung}: Jede gestartete Update-Sequenz erreicht
    nach endlich vielen Schritten entweder den Zustand
    \texttt{UPDATE\_COMPLETE} oder den Zustand \texttt{ROLLBACK\_DONE},
    auch unter der Annahme beliebiger, wiederholter Fehler.
  \item \textbf{Konsistenz}: Nach Abschluss eines Rollbacks weisen alle
    \glspl{psg} exakt die Software-Version auf, die vor dem
    Update-Versuch installiert war (keine Teilrollbacks).
  \item \textbf{Keine Race Conditions}: Das Mutual-Exclusion-Protokoll
    für den Zugriff auf den ZSG-Flash-Speicher verhindert korrumpierte
    Schreiboperationen bei gleichzeitigen Installationsanfragen.
\end{itemize}

Ergänzend zur Verhaltensspezifikation sollte das
\textbf{kryptographische Protokoll} des Update-Manifests (Signatur,
Verschlüsselung, Replay-Schutz durch Nonce/Sequenzzähler) mit einem
formalen Werkzeug wie \textit{ProVerif} oder \textit{Tamarin Prover}
auf Angriffssicherheit hin analysiert werden. Dabei sind insbesondere
Man-in-the-Middle-Angriffe auf den Backend-ZSG-Kanal und
Replay-Attacken auf den internen Fahrzeugbus zu modellieren.

\subsubsection{Künstliche Intelligenz und adaptives Update-Scheduling}

Das in dieser Arbeit vorgeschlagene Update-Scheduling basiert auf
regelbasierten, statisch konfigurierten Zeitfenstern. Diese Strategie
ist praktikabel, lässt jedoch erhebliches Optimierungspotenzial
ungenutzt. Zukünftige Forschung sollte untersuchen, inwiefern
\textbf{maschinelle Lernverfahren} eine dynamische, kontextadaptive
Scheduling-Strategie ermöglichen können.

Relevante Eingabevariablen für ein solches lernbasiertes Scheduling-Modell
umfassen:

\begin{itemize}
  \item \textbf{Fahrzeugzustand}: Batteriefüllstand,
    Motortemperatur, aktive Fahrerassistenzsysteme, Parkstatus.
  \item \textbf{Netzwerkverfügbarkeit}: Signalstärke und
    Datenrate der aktuellen Mobilfunkverbindung, prognostizierte
    Verfügbarkeit in den nächsten Stunden anhand von Bewegungsmustern.
  \item \textbf{Nutzungsverhalten}: Historisch gelernte
    Fahrzeug-Standzeiten, typische Abfahrtszeiten, saisonale
    Muster.
  \item \textbf{Update-Dringlichkeit}: Kritikalitätsklassifizierung
    der vorliegenden Software-Pakete (Security Patch vs.\ Komfort-Update)
    aus dem Software Update Manifest.
  \item \textbf{Flottendaten}: Aggregierte, datenschutzkonforme
    Signale über Update-Erfolgsraten und Rollback-Häufungen für
    bestimmte Software-Versionen auf Flottenebene.
\end{itemize}

Ein vielversprechender Ansatz ist die Formulierung des Scheduling-Problems
als \textbf{Markov Decision Process} (MDP), lösbar durch
\textit{Reinforcement Learning}. Der Agent lernt hierbei, den optimalen
Zeitpunkt für den Beginn eines Update-Vorgangs zu bestimmen, so dass
sowohl die Nutzerunterbrechung minimiert als auch kritische
Sicherheits-Updates innerhalb vorgegebener Fristen eingespielt werden.
Herausforderungen bestehen dabei in der Sicherstellung von
Safety-Garantien des RL-Agenten (keine Aktualisierung von
\gls{asil}-D-Systemen während der Fahrt), die durch formale
\textit{Shields} oder \textit{Safe RL}-Verfahren adressiert werden können.

Auf \textbf{Flottenebene} eröffnet ein solches System die Möglichkeit
einer koordinierten, gestaffelten Einführung neuer Software-Versionen
(\textit{Staged Rollout}): Der Backend-Server kann zunächst eine
Pilotgruppe (z.\,B.\ 1\,\% der Flotte) mit einer neuen Version
aktualisieren, Telemetriedaten auswerten und erst bei ausbleibendem
Fehleranstieg den Rollout auf die Gesamtflotte ausweiten – ein
Deployment-Modell, das aus der Cloud-Software-Entwicklung bekannt ist
und im Automotive-Kontext noch weitgehend unerforscht bleibt.

\subsubsection{Redundanz, Hochverfügbarkeit und Safety-Architektur}

Die vorliegende Arbeit identifiziert das ZSG explizit als
\emph{Single Point of Failure}. Für Fahrzeugfunktionen höchster
Kritikalität – insbesondere für zukünftige vollautonome Fahrzeuge
(SAE Level\,4/5) – ist eine redundante ZSG-Architektur unumgänglich.
Folgende Topologien sind zu untersuchen:

\begin{itemize}
  \item \textbf{Aktiv-Passiv-Redundanz}: Ein primäres ZSG führt sämtliche
    Update-Koordina\-tions\-aufgaben durch; ein sekundäres ZSG synchronisiert
    fortlaufend seinen Zustand und übernimmt bei Ausfall des Primärsystems
    innerhalb eines definierten Zeitfensters (\textit{Failover}).
    Zu klären ist das Protokoll für die Zustandssynchronisierung und
    die Behandlung von Paket-Downloads, die zum Zeitpunkt des Ausfalls
    noch nicht abgeschlossen waren.
  \item \textbf{Aktiv-Aktiv-Redundanz mit Lastverteilung}: Beide ZSGs
    sind parallel aktiv und verteilen die Update-Last auf unterschiedliche
    PSG-Gruppen. Dies erhöht nicht nur die Ausfallsicherheit, sondern
    verkürzt auch die Gesamtdauer von Fahrzeug-weiten Updates.
  \item \textbf{ZSG-Integration in die Zonensteuergeräte}: Als
    Alternative zu einem dedizierten ZSG-Board könnte die
    ZSG-Funktionalität in eines der Zonensteuergeräte oder in den
    zentralen HPC migriert werden, wobei eine klare
    Software-Partitionierung mittels Hypervisor die Isolation
    sicherstellt. Diese Option ist besonders relevant für
    kostenoptimierte Fahrzeugplattformen.
\end{itemize}

Aus Sicht der \textbf{funktionalen Sicherheit} (ISO\,26262) wäre eine
vollständige \gls{asil}-Dekomposition der ZSG-Softwarekomponenten
(\texttt{OTA\_Coordinator}, \texttt{Package\_Vault}, \texttt{Cryp\newline to\_Verifier}) durchzuführen, um die für jede Komponente
geltenden ASIL-Anforderungen abzuleiten. Besondere Aufmerksamkeit verdient
hierbei die Frage, ob der \texttt{OTA\_Coordinator} selbst
\gls{asil}-B oder -D erreichen muss, wenn er die Aktualisierung
von \gls{asil}-D-Subsystemen steuert.

\subsubsection{Post-Quanten-Kryptographie und langfristige Sicherheit}

Die in dieser Arbeit eingesetzten kryptographischen Verfahren –
ECDSA mit P-256 für Signaturen, AES-GCM für die symmetrische
Verschlüsselung und TLS\,1.3 für den Backend-Kanal – entsprechen dem
aktuellen Stand der Technik und sind für die kommenden Jahre als sicher
einzustufen. Allerdings werden \textbf{Quanten-Computer} mittelfristig
die Sicherheit elliptischer-Kurven-basierter Verfahren gefährden.
Das US-amerikanische \textit{National Institute of Standards and
Technology} (NIST) hat 2024 mit \textit{ML-KEM} (CRYSTALS-Kyber)
und \textit{ML-DSA} (CRYSTALS-Dilithium) die ersten
\textbf{Post-Quanten-Kryptographie (PQC)}-Standards verabschiedet.

Zukünftige Arbeiten sollten untersuchen, wie diese Verfahren in das
beschriebene Protokoll integriert werden können:
\begin{itemize}
  \item \textbf{Hybride Signaturverfahren}: Parallel zur bestehenden
    ECDSA-Signatur wird eine ML-DSA-Signatur über das Software Update
    Manifest geführt. Erst wenn beide Signaturen verifiziert sind, gilt
    ein Paket als authentisch. Dieser hybride Ansatz schützt sowohl
    gegen klassische als auch gegen Quanten-Angriffe,
    ohne das bestehende Ökosystem sofort ablösen zu müssen.
  \item \textbf{Key-Encapsulation auf dem Backend-Kanal}:
    Der TLS-Handshake zwischen Backend und ZSG wird um einen
    ML-KEM-basierten Key-Exchange-Mechanismus erweitert,
    so dass der Sitzungsschlüssel auch gegen Quanten-Angriffe
    geschützt ist (\textit{Harvest-now, decrypt-later}-Szenario).
  \item \textbf{HSM-Anforderungen}: Da PQC-Operationen erheblich
    rechenintensiver sind als ECC-Operationen, müssen zukünftige
    \glspl{hsm} explizit für ML-DSA/ML-KEM ausgelegt sein.
    Die Auswirkungen auf die in \cref{sec:evaluation} ermittelten
    Verifikationszeiten sind zu quantifizieren.
\end{itemize}

Parallel dazu bleibt die \textbf{Langzeitsicherheit der
PKI-Infrastruktur} eine offene Herausforderung. Fahrzeuge haben
eine Lebensdauer von 15 bis 20 Jahren; die im ZSG hinterlegten
Root-Zertifikate müssen über diesen gesamten Zeitraum hinweg
rotiert und erneuert werden können. Ein standardisiertes Protokoll
für die sichere \textbf{In-Field-Zertifikatsrotation} –
analog zum ACME-Protokoll in der Webinfrastruktur – ist für
den Automotive-Kontext bislang nicht etabliert und bietet
daher erheblichen Forschungsbedarf.

\subsubsection{Feature-on-Demand, Geschäftsmodelle und regulatorische Einbettung}

Das in dieser Arbeit beschriebene System ist primär auf die technische
Zuverlässigkeit von Software-Updates ausgerichtet. Eine zunehmend
wichtige Dimension betrifft jedoch die \textbf{geschäftlichen und
regulatorischen Rahmenbedingungen}, die das OTA-System direkt beeinflussen.

\paragraph{Feature-on-Demand (FoD).}
Moderne Fahrzeughersteller monetarisieren Software durch das nachträgliche
Freischalten von Fahrzeugfunktionen gegen Entgelt – etwa Sitzheizung,
Beschleunigungsprofile oder Fahrerassistenzsysteme. Das Software Update
Manifest (\gls{sum}) muss hierfür um \textbf{Lizenzierungsmetadaten}
erweitert werden: Jedes Software-Paket wird mit einer
Lizenz-ID, einem öffentlichen Aktivierungsschlüssel und
einer Gültigkeitsdauer versehen. Das ZSG prüft vor der Aktivierung
des Pakets auf dem PSG, ob eine gültige Lizenz im ZSG-internen
Lizenzstore hinterlegt ist. Zukünftige Arbeiten sollten ein
vollständiges \textbf{Lizenzprotokoll} spezifizieren, das
insbesondere folgende Fälle abdeckt: Fahrzeugverkauf mit
Lizenzübertragung, temporäre Aktivierungen (Mietfahrzeuge),
und Lizenzrücknahme bei Zahlungsausfall.

\paragraph{Regulatorische Anforderungen.}
Die Einführung von \textbf{UN-Regelung Nr.\,156} (Cybersecurity
Management Systems, CSMS) und \textbf{UN-Regelung Nr.\,155}
(Software Update Management Systems, SUMS) im Jahr 2021 hat erstmals
verbindliche Anforderungen an OTA-Update-Systeme definiert, die von
Fahrzeugherstellern in der EU, Japan und Südkorea eingehalten werden
müssen~\cite{iso21434}. Die in dieser Arbeit vorgeschlagene Architektur
erfüllt die strukturellen Anforderungen dieser Regelwerke, jedoch fehlt
eine \textbf{systematische Konformitätsprüfung}. Zukünftige Arbeiten
sollten jeden SUMS-Anforderungspunkt der UN-R\,156 auf die vorgeschlagenen
Architekturkomponenten mappen und etwaige Lücken schließen.

Darüber hinaus stellt die \textbf{EU-Produkthaftungsrichtlinie}
(überarbeitete Fassung 2024) Hersteller vor die Frage,
inwiefern ein fehlgeschlagenes OTA-Update als Produktfehler gewertet
werden kann. Die technische Dokumentation des Update-Prozesses –
insbesondere der lückenlose Nachweis von kryptographischer Integrität,
Rollback-Kompletheit und Nutzereinwilligung – gewinnt damit auch
rechtliche Relevanz und sollte als integraler Bestandteil des
OTA-Systems konzipiert werden.

\paragraph{AUTOSAR-Standardisierungsprozess.}
Die vorgeschlagene Erweiterung ist als offen zugängliche
Spezifikation konzipiert. Der nächste konkrete Schritt wäre
die Einbringung eines \textbf{Change Requests} beim
AUTOSAR-Konsortium, insbesondere zielend auf
die Erweiterung von \texttt{AUTOSAR\_SWS\_UCMMaster} um die hier
definierten Schnittstellen \texttt{Ucm\_ParkPackage},
\texttt{Ucm\_ScheduleUpdate} und \texttt{Ucm\_RollbackAll}.
Parallel dazu bietet sich die Formalisierung des
Software Update Manifest als eigenständiges
AUTOSAR-Informationsmodell (\texttt{AUTOSAR\_SRS\_UpdateManifest})
an. Eine Zusammenarbeit mit dem AUTOSAR Working Package
\textit{WP Adaptive Platform Security} und dem
\textit{WP Update and Configuration Management} wäre hierfür
der geeignete organisatorische Rahmen.

\paragraph{Containerisierung und Cloud-native Deployment.}
Längerfristig zeichnet sich im SDV ein Trend zu
\textbf{containerisierten Software-Deployments} ab: AUTOSAR Adaptive
Dienste werden als OCI-konforme Container-Images ausgeliefert,
die auf einem POSIX-betriebssystem mit Container-Laufzeit
(z.\,B.\ \textit{Podman} oder ein automotive-spezifischer
Container-Runtime wie \textit{SOAFD}) ausgeführt werden.
Dies vereinfacht das Update einzelner Dienste erheblich,
da statt eines monolithischen Firmware-Images nur ein
diff-komprimiertes Container-Layer übertragen werden muss.
Die in dieser Arbeit beschriebenen Mechanismen für
Signaturprüfung, Staging und Rollback sind auf containerisierte
Deployments prinzipiell übertragbar; die konkreten
Anpassungen des Protokolls und des Manifest-Formats
sind jedoch Gegenstand zukünftiger Untersuchungen.

\vspace{0.5em}
Zusammenfassend lässt sich festhalten, dass die vorliegende Arbeit
den konzeptionellen Raum eines AUTOSAR-konformen, fahrzeugweiten
OTA-Update-Systems vollständig absteckt, während der
Forschungsbedarf in den Dimensionen experimenteller
Validierung, formaler Absicherung, adaptiver Intelligenz,
kryptographischer Zukunftssicherheit sowie
geschäftlicher und regulatorischer Integration
noch erheblich ist. Der durch diese Arbeit gelegte Grundstein
macht deutlich, dass ein standardisierter, offener Ansatz
gegenüber proprietären Lösungen entscheidende Vorteile
in Bezug auf Interoperabilität, Wiederverwendbarkeit
und Auditierbarkeit bietet – und damit eine
lohnende Investition für das gesamte
AUTOSAR-Konsortium darstellt.

%% Literaturverzeichnis

\addcontentsline{toc}{section}{Literaturverzeichnis}
\begin{thebibliography}{14}

\bibitem{broy2006challenges}
M.~Broy, „Challenges in Automotive Software Engineering," in
\textit{Proc.\ 28th Int.\ Conf.\ Software Engineering (ICSE)},
New York, NY, USA: ACM, 2006, S.~33--42.
\href{https://doi.org/10.1145/1134285.1134292}{doi:10.1145/1134285.1134292}

\bibitem{nilsson2008key}
D.~K.~Nilsson und U.~E.~Larson, „Secure Firmware Updates over the Air in
Intelligent Vehicles," in \textit{Proc.\ ICC Workshops 2008},
IEEE, 2008, S.~380--384.
\href{https://doi.org/10.1109/ICCW.2008.77}{doi:10.1109/ICCW.2008.77}

\bibitem{autosar2022ucm}
AUTOSAR Consortium, \textit{Specification of Update and Configuration
Management}, AUTOSAR Technical Specification
AUTOSAR\_SWS\_UpdateAndConfigurationManagement, 2022.
\url{https://www.autosar.org/}

\bibitem{autosar2022ap}
AUTOSAR Consortium, \textit{Explanation of Adaptive Platform Design},
AUTOSAR Technical Report AUTOSAR\_EXP\_AdaptivePlatformDesign, 2022.
\url{https://www.autosar.org/}

\bibitem{autosar2022cp}
AUTOSAR Consortium, \textit{Layered Software Architecture},
AUTOSAR Technical Specification
AUTOSAR\_EXP\_LayeredSoftwareArchitecture, 2022.
\url{https://www.autosar.org/}

\bibitem{autosar2022secoc}
AUTOSAR Consortium, \textit{Specification of Secure Onboard Communication},
AUTOSAR Technical Specification
AUTOSAR\_SWS\_SecureOnboardCommunication, 2022.
\url{https://www.autosar.org/}

\bibitem{iso21434}
ISO, \textit{Road Vehicles -- Cybersecurity Engineering},
ISO/SAE~21434:2021, International Organization for Standardization, 2021.

\bibitem{iso26262}
ISO, \textit{Road Vehicles -- Functional Safety},
ISO~26262:2018, International Organization for Standardization, 2018.

\bibitem{iso14229}
ISO, \textit{Road Vehicles -- Unified Diagnostic Services (UDS)},
ISO~14229-1:2020, International Organization for Standardization, 2020.

\bibitem{halder2020secure}
S.~Halder, A.~Ghosal und M.~Conti, „Secure OTA Software Updates in
Connected Vehicles: A Survey," \textit{Computer Networks},
Bd.~178, S.~107343, 2020.
\href{https://doi.org/10.1016/j.comnet.2020.107343}{doi:10.1016/j.comnet.2020.107343}

\bibitem{zelle2020security}
D.~Zelle, T.~Lauser und C.~Krau{\ss}, „Analyzing and Securing SOME/IP
Automotive Services with Formal and Semi-Formal Methods," in
\textit{Proc.\ 17th Int.\ Conf.\ Security and Cryptography (SECRYPT)},
SCITEPRESS, 2020, S.~24--35.
\href{https://doi.org/10.5220/0009833600240035}{doi:10.5220/0009833600240035}

\bibitem{mundhenk2017security}
P.~Mundhenk et al., „Security in Automotive Networks: Lightweight
Authentication and Authorization," \textit{ACM Trans.\ Design Autom.\
Electronic Systems}, Bd.~22, Nr.~2, S.~25:1--25:27, 2017.
\href{https://doi.org/10.1145/2998388}{doi:10.1145/2998388}

\bibitem{jo2021ovota}
S.~Jo, J.~Park, Y.~Shin und Y.~Kwon, „OVOTA: Toward Secure OTA Update
on Autonomous Vehicles," in \textit{Proc.\ AutoSec 2021},
ACM, 2021.
\href{https://doi.org/10.1145/3459086.3459640}{doi:10.1145/3459086.3459640}

\bibitem{unece2021wp29}
UNECE, \textit{Regulation No.~156 -- Uniform provisions concerning the
approval of vehicles with regards to software update and software updates
management system}, UN/ECE/TRANS/505/Rev.3/Add.155, 2021.

\bibitem{kugele2018sdv}
S.~Kugele, D.~Hettler und J.~Peter, „Data-Centric Communication and
Containerization for Future Automotive Software Architectures," in
\textit{Proc.\ IEEE Int.\ Conf.\ Service-Oriented Computing and Applications
(SOCA)}, IEEE, 2018, S.~65--73.
\href{https://doi.org/10.1109/SOCA.2018.00018}{doi:10.1109/SOCA.2018.00018}

\bibitem{schneider2026ladekomm}
T.~Schneider und M.~Vick, „Ladekommunikation meets SDV -- Ladecontroller in
der Zonenarchitektur," \textit{Elektronikpraxis}, Nr.~1, S.~60--62, Januar 2026.
Vector Informatik GmbH.

\bibitem{iso15118}
ISO, \textit{Road Vehicles -- Vehicle to Grid Communication Interface},
ISO~15118-20:2022, International Organization for Standardization, 2022.

\bibitem{autosar2022ucm_master}
AUTOSAR Consortium, \textit{Specification of UCM Master},
AUTOSAR Technical Specification
AUTOSAR\_SWS\_UpdateConfigurationManagementMaster, 2022.
\url{https://www.autosar.org/}

\end{thebibliography}


%% Abkürzungsverzeichnis

\printglossary[type=\acronymtype,title=Abkürzungsverzeichnis]

\end{document}

